\pagewiselinenumbers
\setlength{\footmarkwidth}{1.8em}
\label{edition01}
%\tstl{hari}
\var{$\ast$}{\cm Before the beginning: \dn [ह]रि श्रीगणपतये नमः [अविघ्न]मस्तु॥ \Tm, \dn पतञ्जलये नमः \LMA}%\donote{hari}{Cf.\ the beginnig of the manuscript of the \Svaditamkarani\ by \Paramesvara~I on the \Nyayakanika: \dn हरिः श्रीगणपतये नमः गणपतिरभिमतमनिशं दिशतु\cm\ \citep[59]{vidhiv:stern} with \Tm. In that sense, the statement made by Sarma 1991 about the ``generic statement'' in the manuscript of Yuktibhāṣā should also be looked at.}

\begin{verse}
\mds%\com{zero}{यस्मिन्न स्तः \ldots\ प्रणमामि}
\llbl{p01_1}\labelh{FVFPEDN}%
\tstwll{FRSTVRS1}{यस्मिन्न स्तः \ldots\ कैटभशत्रुम्प्रणमामि}%
\tstwll{YTASTM1}{यस्मिन्न स्तः \ldots\ निखिलानाम्}%
\varl{2}य\mms स्मिन्न\mde\ स्तः\mme\ कर्म्मविपाकौ %$\dagger$
यत \mms आस्तां%$\dagger$
\donote{2}{%
  \vnote{L\-\Mpc A\-\Me}{%
    \rdg{[य\mist{3}कर्म्मवि\ccs वा\cce]पा[कौ यत\mist{3}]}{\Tm},
    \rdg{\mist{3}स्तः कर्मविपाको यत आस्तां}{\Td},
    \rdg{यस्मिन्नन्तः कर्मविपाकेन यत आस्तां}{\Mac}}} 
\\ \hspace{60pt} %
% \vrt{\XeTeXglyph 175 \XeTeXglyph 349\mme शा य\mms स्मै\mme\ \mds ना\mms लमलंघ्या\mde\ निखिला\mme नाम्।}%
\vrt{\kle\mme शा य\mms स्मै\mme\ \mds ना\mms लमलंघ्या\mde\ निखिला\mme नाम्।}%
  {\LMA\-\Me}{%
    \rdg{[\mist{2}ा य\mist{7}ा]नाम्}{\Tm}, 
    \rdg{\kle शा यस्मै \mist{5} नखिलानाम्}{\Td}}%
\donote{YTASTM1}{See \rYS{1.24}{1.24}: {\dn \kle शकर्मविपाकाशयैरपरामृष्टः \barfu पुरुषविशेष \barfu \barfu ईश्वरः॥}; Cf.\ \SVSAP{069--072}{kk.\ 69--72}: {\dn सर्वेषां तु फलापेतं न स्थानमुपपद्यते। न \barfu चाप्यनुपभोगो ऽसौ कस्यचित्कर्मणः फलम्॥७०॥ अशेषकर्मनाशे वा पुनःसृष्टिर्न युज्यते। कर्मणां \barfu वाप्यभिव्यक्तौ किं निमित्तं \barfu तदा भवेत्॥७१॥ ईश्वरेच्छा यदीष्येत सैव स्याल्लोककारणम्। ईश्वरेच्छावशित्वे हि निष्फला \barfu कर्मकल्पना॥७२॥}, 75: {\dn कस्यचिद्धेतुमात्रत्वं यद्यधिष्ठातृतेष्यते। कर्मभिः सर्वजीवानां तत्सिद्धेः सिद्धसाधनम्॥}.}%
\\
\tstwll{KLYT1}{नावच्छिन्नः \barfu कालदृशा \barfu यः \barfu कलयन्त्या}नावच्छिन्नः %
\vrt{कालदृशा}{LMA\Me}{%
    \rdg{कालदिशा}{T}%
	%\rdg{कालदृशा}{LMA\Me}%
} %
यः कलयन्त्या%
\donote{KLYT1}{See \rYS{1.26}{1.26}: {\dn स पूर्वेषाम\-पि गुरुः कालेनानवच्छेदात्॥}.  Cf.\ \rBhG{10.30b}{10.30b}: {\dn कालः कलयतामहम्} and \rnnBhGBh{10.30}{10.30} on it: {\dn कालः कलयतां कलनं गणनं कुर्वतामहम्}.} %
\\ \hspace{60pt} \vrt{लोकेशन्तं}{TL}{\rdg{लोकेशस्तं}{MA\Me}}
%\resume{one}
कैटभशत्रु\mts म्प्रणमामि॥१॥%
\donote{FRSTVRS1}{Cf.\ the salutatory stanza of the \NKn: %
    {\dn परामृष्टः \kle शैः कथमपि न यो जातु भगवान्न धर्माधर्माभ्यां त्रिभिरपि विपाकैर्न च \barfu तयोः। परं वाचां तत्त्वं यमधिगमयत्योमिति पदं  न\-म\-स्या\-मो विष्णुं तममरगुरूणामपि गुरुम्॥}. %
	See also \Paramesvara's commentaries (the \Jusadhvamkarani\ and the \Svaditamkarani, especially the latter) on the \NKn\ where the YVi is referred to \citep[106--25]{vidhiv:stern}. %
	Cf.\ also \index{Halbfass, Wilhelm}\citet[207]{halbfass:sankara}.%
}%\donote{zero}{This verse is in the Mattamayūra meter.}
%\bigskip
% \newpage
%\com{versetwo}{यस्सर्व्ववित् \ldots गुरोरपि}
% \tstwll{SRVVVS1}{यस्सव्ववित् सर्व्वविभूतिशक्तिः}यस्स\mme र्व्व\mde वि\mts त्\mte\  सर्व्वविभूतिशक्ति%
\\ \tstwll{SRVVVS1}{यस्सव्ववित् सर्व्वविभूतिश\kti ः}\llbl{p01_2}यस्स\mme र्व्व\mde वि\mts त्\mte\  सर्व्वविभूतिश\kti%
\donote{SRVVVS1}{See \pageref{ANUMANAONE},\lineref{ANUMANAONE}--\pageref{thrprfe},\lineref{thrprfe}.}%
\tstwll{VIHIND1}{विहीनदोषोपहितक्रियाफलः}%
र्व्विहीन%
\vrtmm{दोषोपहित}{\Tmpc LMA\Me}{%
    \rdg{॰दोषो\ins प\ine हित॰}{\Tm}, 
	\rdg{॰दोषो ऽपि हि तत्}{\Td}}%
क्रिया%
\vrtms{फलः}{M(em.?)A\Me}{%
    \rdg{॰फलम्}{TL}}%
\donote{VIHIND1}{See %
    \pageref{PASRR25B},\lineref{PASRR25B}--\pageref{PASRR25E},\lineref{PASRR25E};  %
	Cf.\ \pageref{DSSR25B},\lineref{DSSR25B}--\pageref{SRRTV25E},\lineref{SRRTV25E}.}।%
%{conj.}{%
%    \rdg{॰र्व्वीहीनदोषो\ins प\ine हितक्रियाफलम्}{\Tm},
%    \rdg{॰र्विहीनदोषो ऽपि हि तत्क्रियाफलम्}{\Td},
%    \rdg{॰र्विहीनदोषोपहितक्रियाफलम्}{L},
%    \rdg{र्वि\vv{॰: वि॰}{E}हीनदोषोपहितक्रियाफलः}{MA\Me}}% 
 \\
विश्वोद्भवा\vrtmm{न्तस्थिति}{\Tm LMA\Me}{%
  \rdg{॰न्तः स्थिति॰}{\Td}%
%  \rdg{(विश्वोद्भवान्तस्थिति॰)}{L}%
}हेतुरीशो \tstwll{GRGR1}{नमो ऽस्तु तस्मै गुरवे गुरोरपि}\vrt{नमो ऽस्तु}{T\-L\vv{नमो स्तु}{\Tm L}\-A(em.?)\-\Me(em.)}{\rdg{नम स्तु}{M}} तस्मै %\resume{v2}
\vrt{गुरवे}{\Td LMA\Me}{\rdg{गुरवै}{\Tm}} गुरोरपि\donote{GRGR1}{See \rYS{1.26}{1.26}: {\dn  स पूर्वेषाम\-पि गुरुः कालेना\-नव\-च्छेदात्}.} ॥२॥%\donote{v2}{For b-p\=ada, See \rYS{1.24} and YBh 1.24 ({\dn स हितफलस्य भोक्ता}); d-p\=ada, see \rYS{1.26}.}%\donote{versetwo}{\come{versetwo}}
\end{verse}
\newpage

\llbl{p01_3}\mll{ys1}\str{\vrt{अथ \barfu योगानुशासनम्}{TLMA}{om.\ \Me}॥१॥}\donote{ys1}{{\dn अथ योगानुशासनम्॥}}

% \bigskip
%\bhsy{अथेत्ययमधिकारार्थः। योगानुशासनं शास्त्रमधिकृतं वेदितव्यम्}
\labelh{FIRSTSTC}%
\llbl{p01_4}\tstwll{FSTSTC1}{अथेत्यादि पातञ्जलयोगशास्त्रम्; तस्य विवरणमारभ्यते}%\begin{uncertain}%
अथेत्यादि पातञ्जलयोग%
\vrtmm{शास्त्र}{TLMA}{%
    \rdg{॰शास्त्र\eil सूत्रभाष्य\eir ॰}{\Me}}%
\vrtms{म्; तस्य}{conj.}{%
    \rdg{\om}{\all}} %
विवरणमारभ्यते%\end{uncertain}%
।%
\donote{FSTSTC1}{\labelh{OPENG}%
    Cf.\ the openings of the \rnnBhGBh{00intro}{intro.}:
    {\dn तदिदं गीताशास्त्रं समस्तवेदार्थसारसंग्रहभूतं दुर्विज्ञेयार्थं तदर्थाविष्करणायानेकैर्विवृतपदपदार्थवाक्यवाक्यार्थन्यायमप्यत्यन्तविरुद्धानेकार्थत्वेन लौकिकैर्गृह्यमाणमुपलभ्याहं विवेकतो ऽर्थनिर्धारणार्थं संक्षेपतो विवरणं करिष्यामि।} \citep[5--6]{bhgbh}; %
    \rnnChUBh{00intro.}{intro.}: %
    {\dn ओमित्येतदक्षरमित्याद्यष्टा(॰दिरष्टा॰} or {\dn ॰द्याष्टा॰?)ध्यायी छान्दोग्योपनिषत्। तस्याः सङ्क्षेपतो ऽर्थजिज्ञासुभ्यः ऋजु विवरणमल्पग्रन्थमिदमारभ्यते।} \citep[1]{upa2}; %
    \BAUBh: %
    {\dn `उषा वा अश्वस्य' इत्येवमाद्या वाजसनेयिब्राह्मणोपनिषत्। तस्या इयमल्पग्रन्था वृत्तिरारभ्यते।} \citep[1]{bau}; %
    \KeUBh:
    {\dn `केनेषितमि'त्याद्योपनिषत्परब्रह्मविषया वक्तव्येति \barfu नवमस्याध्यायस्यारम्भः।} \citep[17]{upa1};
    \KaUBh:
    {\dn अथ काठकोपनिषद्वल्लीनां सुखेनार्थप्रबोधनार्थमल्पग्रन्था वृत्तिरारभ्यते।} \citep[55]{upa1};
    \MuUBh:
    {\dn `ओं \barfu ब्रह्मा देवानामि'त्याद्याथर्वणोपनिषत्।} \citep[127]{upa1};
    \MaUBh:
    {\dn `ओमित्येतदक्षरमिदं सर्वं तस्योपव्याख्यानम्।' वेदान्तार्थसारसंग्रहभूतमिदं प्रकरणचतुष्टय`मोमित्येतदक्षरमि'त्याद्यारभ्यते।} \citep[212]{upa1}} %
तत्रानाख्यातसम्बन्धप्रयोजनन्न %\tstwll{PPNN}{पुरुषप्रवृत्तिनिवृत्तिभ्याम्}
पुरुष\-प्रवृत्ति\-निवृत्ति\-भ्या\-%\donote{PPNN}{blah}
म्पर्या\-प्नोतीति सूत्रकाराभिप्रेते पुरुष\-प्रवृत्ति%
\vrtmm{निवृत्ति\-नि\-मि\-त्त}{\conj}{%
    \rdg{॰निमित्त॰}{\all}}%
\vrtms{भूते}{TMA\Me}{%
\rdg{॰भूतैः}{L}} %
\tstwll{SBPY}{सम्बन्धप्रयोजने प्रकटीक्रियेते}सम्बन्धप्रयोजने पूर्व्वम्प्रकटीक्रियेते॥%
\donote{SBPY}{Cf.\ \SVPRA{011--25}{kk.\ 11--25}: {\dn    
	`अथातो धर्मजिज्ञासा' सूत्रमाद्यमिदं कृतम्।     धर्माख्यं विषयं वक्तुं मीमांसायाः प्रयोजनम्॥११॥    
	सर्वस्यैव हि शास्त्रस्य कर्मणो वापि कस्यचित्।     यावत् प्रयोजनं नोक्तं तावत् तत् केन गृह्यते॥१२॥     
	मीमांसाख्या तु विद्येयं \barfu बहुविद्यान्तराश्रिता।     न शुश्रूषयितुं शक्या प्रागनुक्त्वा प्रयोजनम्॥१३॥     
	विद्यान्तरेषु नाप्येतद्यद्यभीष्टं प्रयोजनम्।     अनर्थप्रापणं तावत् तेभ्यो नाशङ्क्यते क्वचित्॥१४॥     
	मीमांसायां त्विहाज्ञाते दुर्ज्ञाते \barfu वाविवेकतः।     न्यायमार्गे महान् दोष इति \barfu यत्नोपचर्यता॥१५॥     
	तस्मात्प्रयोजनं पूर्वमुक्तं सूत्रकृता स्वयम्।     यत्स्वेनोक्तं वदेयुस्तद्भाष्यकारादयः कथम्॥१६॥     
	सिद्धार्थं \barfu ज्ञातसम्बन्धं \barfu श्रोतुं श्रोता प्रवर्तते।     शास्त्रादौ तेन \barfu वक्तव्यः \barfu सम्बन्धः \barfu सप्रयोजनः॥१७॥     %
	शास्त्रं प्रयोजनं चैव सम्बन्धस्याश्रयावुभौ।     तदुक्त्यन्तर्गतस्तस्माद्भिन्नो नोक्तः प्रयोजनात्॥१८॥     
	सिद्धिः श्रोतृप्रवृत्तीनां \barfu सम्बन्धकथनाद्यतः।     तस्मात्सर्वेषु शास्त्रेषु सम्बन्धः पूर्वमुच्यते॥१९॥     
	यावत्प्रयोजनेनास्य सम्बन्धो नाभिधीयते।     \barfu असम्बद्धप्रलापित्वाद्भवेत्तावदसंगतिः॥२०॥     
	इह त्वाक्षिप्य सम्बन्धं भाष्य एवाभिधास्यते।     धर्मप्रसिद्ध्यसिद्धिभ्यां तस्मान्नान्यो ऽभिधीयते॥२१॥     
	न चाप्यत्राथशब्देन शास्त्रसम्बन्ध उच्यते।     सम्बन्धक्रिययोर्ह्येष ब्रूते शास्त्राच्च ते पृथक्॥२२॥     
	यो ऽप्ययं शास्त्रसम्बन्धो वर्ण्यते \barfu कैश्चिदादितः।     क्रियान्तर्यरूपो वा गुरुपर्वक्रमो ऽपि \barfu वा॥२३॥     
	तदतद्भावयोस्तस्य विशेषो नोपलभ्यते।     श्रोतुर्विधौ निषेधे वा ज्ञाने वा शास्त्रगोचरे॥२४॥     
	\barfu तस्माद्व्याख्याङ्गमिच्छद्भिः \barfu सहेतुः \barfu सप्रयोजनः।     %
	शास्त्रावतारसम्बन्धो वाच्यो नान्यस्तु \barfu निष्फलः॥२५॥}; %
    \rnnNBh{1.1.1}/\rnnNV{1.1.1}\ 1.1.1, and the opening discussion of the \VMBh.  Cf.\ also \dnfiller the introduction to the \BhGBh: {\dn विशिष्ट\-प्रयोजन\-संबन्धा\-भिधेय\-वद्गीता\-शास्त्रं यत\-स्तद\-र्थ\-विज्ञानेन समस्त\-पुरुषा\-र्थ\-सिद्धि\-रत\-स्तद्वि\-वरणे यत्नः क्रियते \barfu मया।} \citep[7]{bhgbh}; \rMaUBh{1.01}{1.1}: {\dn वेदान्ता\-र्थ\-सार\-संग्रह\-भूत\-मिदं प्रकरण\-चतुष्टयमो\-मि\-त्ये\-तद\-क्षर\-मि\-त्या\-द्या\-रभ्यते। अत एव न पृथक्सम्बन्धाभिधेयप्रयोजनानि वक्तव्यानि। यान्येव तु वेदान्ते सम्बन्धा\-भिधय\-प्र\-यो\-जना\-नि तान्ये\-वे\-ह भवितुम\-र्ह\-न्ति। तथा\-पि प्र\-करण\-व्याचिख्या\-सुना संक्षेपतो वक्तव्यानि। तत्र प्रयोजन\-वत्साधना\-भिव्यञ्जकत्वेनाभिधेयसम्बद्धं शास्त्रं पारम्पर्येण विशिष्टसम्बन्धाभिधेयप्रयोजनवद्भवति। किं पुन\-स्त\-त्प्रयोजन\-मित्युच्यते। रोगार्तस्येव \barfu रोग\-नि\-वृत्तौ \barfu स्व\-स्थ\-ता। तथा दुःखात्मकस्यात्मनो द्वैतप्रपञ्चोपशमे \barfu स्व\-स्थ\-ता अद्वैत\-भावः प्रयोजनम्।} \citep[212--3]{upa1}; \rKaUBh{1.1.01}{1.1.1}: {\dn प्रयोजनं चास्या उपनिषद \barfu आत्यन्तिकी \barfu संसारनिवृत्तिर्ब्रह्मप्राप्तिलक्षणा। \barfu सम्बन्धश्चैवंभूतप्रयोजनेनोक्तः।} \citep[58]{upa1}; and \rMuUBh{1.1.01}{1.1.1}. Cf.\ also the beginning of the \VMBh\ where the purpose (\prayojana) of the \emph{śabdānuśāsana} and the relationship between the speech (\sabda) and the meaning (\artha) are discussed in succession.}%

% \newpage
\llbl{p01_5}तत्र प्रयोजनन्तावत्---

\tolerance=1000
\emergencystretch=1.5pt

\llbl{p01_6}चिकित्सा\varmm{शास्त्र इव तच्चतुर्व्यूहत्व}{em.,
  {\dn ॰शास्त्रेवतच्चतुर्व्यूहत्व॰} \Tm L,
  \rdg{॰शास्त्रवत्तच्चतुर्व्यूहत्व॰}{\Td},
%  {\dn ॰शास्त्रवतच्चतुर्व्यूहश्च} L,
  {\dn ॰शास्त्रे तच्चतुर्व्यूहश्च} MA,
  {\dn ॰शास्त्रे तच्चतुर्व्यूहत्व॰} \Me%
}प्रदर्शन\varms{द्वारेण}{T\Me,
 {\dn ॰चारेण} LMA%
} व्याख्यातम्। \barfu तद्यथा %
\tst{चिकित्साशास्त्रञ्च\-तु\-र्व्यू\-ह\-म्---रोगो रोगहेतुरारोग्यम्भैषज्यमिति}{\rYBh{2.15}{2.15}.}% \com{vdhpnym}{विधिप्रतिषेधनियमद्वारेण}
---विधिप्रतिषेध%
\varms{नियमद्वारेण}{em.,
% Probably I need to explain this. See YVi p. 184.
  {\dn ॰नियमचारेण} TLMA,
  {\dn ॰नियमद्वारेण \eil च तत्\eir } \Me(conj.)} %\donote{vdhpnym}{\come{vdhpnym}}
चतुर्व्यूह%
\varmm{विषय}{\Td LMA\Me, {\dn ॰विषया॰} \Tm}\-%
व्या\-ख्या\-न\-परम्। \var{एव\-मिहापि}{TMA\Me, {\dn एवमिहावि॰} L}  \tstl{ys2.15}\quoteo परिणाम\-\pause{ys2.15}<\selm>\sf{\dn ताप\-}\varmm{\dn संस्कार\-दुःखै}{TL\Mpc A\Me, {\dn ॰दुःखसंस्कारै॰} \Mac}%
र्गुण\-वृत्ति\-विरोधाच्च दुःखमेव सर्व्वं \resume{ys2.15}विवेकिन\quotec%
\donote{ys2.15}{%
	\rYS{2.15}{2.15}.  See also \citet{wezler:yvi2}\index{Wezler, Albrecht}.%
} %
इत्यारभ्य चतुर्व्यूहत्वं शास्त्रस्य प्रदर्शितम्। \tstl{ybh215a}\pause{ybh215a}<{\dn दुःखप्रचुरः \barfu संसारो}\sel{\dn तद्दृशेः कैवल्यम्}>तद्यथा---%
\tst{दुःख\-%
\varms{प्रचुरः \barfu संसारो}{%
	M(em.)A\Me, 
	{\dn ॰प्र\-चुर\-सं\-सा\-रो} TL%
} %
\barfu हेयः}{%
	Cf.\ \rYS{2.16}{2.16}: {\dn हेयं दुःखमनागतम्॥}.%
}; %
\tstl{ys2.17}%
त\-स्या\-%
\pause{ys2.17}<\selm>विद्या\-\mms नि\mme मित्तो \barfu द्रष्टृ%
\varmm{दृश्य}{%
	TMA\Me, 
	{\dn ॰दृश्यं} L%
}%
\resume{ys2.17}%
संयोगो \barfu  हेतुः%
\donote{ys2.17}{%
	Cf.\ \rYS{2.17}{2.17}: {\dn द्रष्टृदृश्ययोः \barfu संयोगो \barfu हेयहेतुः॥}.%
}%
; %
\tst{विवेकख्यातिरविप्लवा \barfu हानोपायः}{Cf.\ \rYS{2.26}{2.26}: {\dn विवेकख्यातिरविप्लवा \barfu हानोपायः॥}.}%
; \barfu विवेक\-ख्यातौ च सत्यामविद्यानिवृत्तिः, %
\tst{तन्निवृत्तावात्यन्तिको द्रष्टृदृश्यसंयोगोपरमो हानम्; तदेव कैवल्यमिति।} %
    {Cf.\ \rYS{2.25}{2.25}: {\dn तदभावात्संयोगाभावो हानम्, तद्दृशेः कैवल्यम्॥}.}%
%
\donote{ybh215a}{%
	Cf.\ \rYBh{2.15}{2.15}: {\dn तत्र दुःखबहुलः संसारो \barfu हेयः। प्रधान\-पुरुष\-योः \barfu संयोगो \barfu हेयहेतुः। संयोगस्यात्यन्तिकी निवृत्तिर्हानम्। हानोपायः सम्यग्दर्शनम्।}. %
	Cf.\ also \rNBh{1.1.1}: {\dn हेयम्, तस्य निर्वर्तकम्, हानमात्यन्तिकम्, तस्योपायो ऽधिगन्तव्य इत्येतानि चत्वार्य्यर्थपदानि \barfu सम्य\-ग्बुद्ध्वा निःश्रेयस\-म\-धि\-गच्छति।} \citep[33]{ns}; %
	\rNV{1.1.1}: {\dn हेय\-हानो\-पाया\-धि\-गन्तव्य\-भेदा\-च्चत्वा\-र्य्यर्थ\-पदा\-नीति। हेयं दुःखं तद्धेतुश्च। दुःखमुक्तं। हेतुरविद्यातृष्णे धर्माधर्माविति। हानं तत्त्वज्ञानम्। तत् पुनर्यथार्थावस्थितपदार्थाधिगतिस्तच्च प्रमाणम्। उपायः शास्त्रं तदप्युक्तम्। \barfu अधिगन्तव्यो \barfu ऽपवर्गः \barfu स \barfu पुनरात्यन्तिको \barfu दुःखाभावः।} \citep[11--2]{ns}.%
} %
आरोग्यस्थानीयकैवल्य%
\varmm{प्रयुक्तत्वा}{\Td LM\Me, {\dn ॰प्रयुत\-त्वा॰} \Tm, {\dn ॰प्रयुक्त्वा॰} A}%
दस्य तदेव कैवल्यम्प्रयोजनम्॥

\tolerance=200
\emergencystretch=0pt

\llbl{p01_7}ननु च हेयतद्धेतू न %
\var{प्रस्तोतव्यौ}{%
	T\Lpc, %
	{\dn प्रस्तोतव्येन} \Lac MA, %
	{\dn प्रस्तोतव्ये} \Me}%
, निष्प्रयोजनत्वात्। हानप्रयुक्तं हि शास्त्रम्। \barfu तदुपायभूता विवेकख्यातिरेव \barfu वक्तव्या। %
\labelh{NVKANTHA}%
\tstl{nv111}\pause{nv111}<न हि कण्टकविद्ध॰\sel चोद्येते>%
न हि कण्टकविद्धचरणतलस्य %
\vrtsm{तदपनयन}{T\-L\-\Mpc\-\Me}{%
	\rdg{तदु\-प\-नयन॰}{\Mac\-A}%
}%
म्मुक्त्वा दुःख%
\varms{तद्धेतू}{%
	TL\Mpc A\Me, %
	{\dn ॰तद्धेतु} \Mac} %
चोद्येते॥%
\donote{nv111}{%
	Cf.~\rBAUBh{4.3.07}{4.3.7}: {\dn कण्टकविद्धस्य हि कण्टकवेधजनितदुःखनिवृत्तिः फलं, न तु कण्टकविद्धमरणे तद्दुःखनिवृत्तिफलस्याश्रय उपपद्यते।}; %
	\rNV{1.1.1}: {\dn श्रेयः पुनः सुखमहितनिवृत्तिश्च, \barfu तच्छ्रेयो भिद्यमानं द्वेधा व्यवतिष्ठते दृष्टा\-दृष्ट\-भेदेन। दृष्टं \barfu सुखमदृष्टमहितनिवृत्तिः। \barfu अहितनिवृत्तिरप्यात्यन्तिकी \barfu अनात्य\-न्ति\-की च। \barfu अना\-त्यन्ति\-की, कण्टकादेर्दुःखसाधनस्य परिहारेण। \barfu आत्यन्तिकी \barfu पुनरेकविंशतिप्रभेदभिन्नदुःखहान्या।}~\citep[5--6]{ns}.%
}

\llbl{p01_8}\mms नैतदेवम्\mme\kern 2pt । हेयतत्कारणापेक्षत्वाद्धानोपायस्य। यावदिदमनेनोपायेन हातव्यं संसारचक्रम्, अस्य \barfu चाविद्या\-%
\varms{निमित्तो}{%
	LMA\Me, 
	{\dn ॰नि\-मि\-तो} \Tm, 
	{\dn ॰निमित्तौ} \Td%
} %
द्रष्टृ\-दृश्य\-संयोगो हेतुरिति नाख्यायते, तावन्नाविद्याप्रतिपक्ष\-भूता विवेक\-ख्याति\-रा\-ख्यातुं शक्यते। हेयतद्धेतु\vrtmm{मतो ऽपि \barfu हानो}{em.}{\rdg{॰मतो हि हानो॰}{\Tmpc\Td\Lpc}, {\dn ॰मतोर्व्विहानो॰} \Tmac, {\dn ॰मतो हि \ccs ना\cce हानो॰} L, \rdg{॰मतो हानो॰}{MA\Me}}पाया\-र्त्थिनो \var{रोगिणो}{TMA\Me, {\dn रोगिणौ} L} भै\-ष\-ज्या\-र्त्थित्वदर्शनात्। न हि रोगं रोग\varmm{हेतुञ्चा}{\Tdpc\Mpc A\Me\vv{\dn ॰हेतुं चा॰}{\Mpc A\Me}, {\dn ॰हेतुस्त्वा॰} \Tm, {\dn ॰हेतुश्चा॰} \Tdac, {\dn ॰हेतुत्वा॰} L, {\dn ॰हेतुं त्वा॰} \Mac}न\varms{पेक्ष्य}{\Td LM\Me, {\dn ॰पेक्ष्या} \Tm, {\dn ॰पयदय॰} A} चिकित्साशास्त्रमुपदिश्यते॥

\llbl{p01_9}\llbl{smb}सम्बन्धो ऽपि---\llbl{smb_end}

\llbl{p01_10}विवेकख्यातेर्हानमेव फलं साध्यम्, \varl{10}हानस्यापि\pause{10}<\sel> विवेकख्यातिरेव \resume{10}साधनमिति\donote{10}{TMA\Me, {\dn हा\-न\-स्या\lost{14}मिति} L} साध्य\-{}साधन\-{}यो\-{}रित\-{}रेतरनियम एव, %
\varsm{नान्यः \barfu सम्ब}{T, {\dn नान्यसम्ब॰}\vv{॰संब॰}{MA\Me}LMA\Me\@}%
न्धः। तद्यथा---\varsm{भै\-ष\-ज्य}{\Td LMA\Me, {\dn भैषज्या॰} \Tm}स्यारोग्यमेव फलम्, आरोग्यस्यापि भैषज्यमेव \barfu साधनमिती\varms{तरेतरनियमः। स च शास्त्रादिति॥}{TMA\Me, {\dn ॰तरेतर\lost{9}स्त्रादिति} L}

\llbl{p01_11}\varsm{तस्मात्सत्सम्ब}{%
	TLMA, {\dn तस्मात्ससंब॰} \Me%
}%
न्धप्रयोजनं योगानुशासनम्॥
  
% \medskip

\llbl{p01_12}ननु च यदि हानम्प्रयोजनम्, तदुपायश्च विवेकख्यातिरिति, तत्र वक्तव्यम`थ विवेकख्यात्यनुशासनमि'ति। किमर्थम`थ योगानुशासनमि'ति सूत्रितम्?

\llbl{p01_13}तदुपायत्वाद्योगस्य। %\com{upayaeva}{उपाय एव वक्तव्य स्यात्}
उपाय एव वक्तव्य स्यात्%
%\donote{upayaeva}{\come{upayaeva}}
---\vrt{उपेय\-प्र\-दर्शने}{em.}{\rdg{उपेयो\-प\-प्र\-दर्शने}{TL\Mpc A\Me}, \rdg{उपेयोप\ins प्र\ine दर्शने}{M}} \varsm{ह्युपेय}{T, {\dn हि उपेय॰} LMA\Me}मुपायम्प्रति साकांक्षमेवेत्युपायः सांगकलापः पुनरपि 
\var{वक्तव्य स्यात्}{%
	TMA\Me\vv{॰व्यः स्या॰}{M\-A\-\Me}, 
	{\dn वक्तव्यः \lost{1}}{\dn ात्} L%
}\fubar । \varsm{तस्मिंस्त्वभि}{TL, {\dn तस्मिंश्चाभि॰} MA\Me}हिते सर्व्वमभिहितमिति॥

\llbl{p01_14}कथन्तदुपायत्वम्?

\llbl{p01_15}यत आह सूत्रकारः---%
\tst{`योगांगानुष्ठानादशुद्धिक्षये ज्ञानदीप्तिरा विवेकख्यातेः'}{%
	\rYS{2.28}{2.28}.%
}%
\varl{iti1}\pause{iti1}<{\dn ॰ख्यातेः'; `निर्व्विचार॰}>\resume{iti1}%
\donote{iti1}{%
	TLMA, 
	{\dn ॰ख्यातेः' इति। `निर्विचार॰} \Me}%
; %
\tst{`निर्व्विचारवैशारद्ये ऽध्यात्मप्रसादः'}{%
	\rYS{1.47}{1.47}.%
}; %
\tst{`ऋतम्भरा तत्र प्रज्ञे'ति}{%
	\rYS{1.48}{1.48}.%
} %
च। तथा भाष्यकारो ऽप्याह---%
\tst{`यस्त्वेकाग्रे चेतसि स भूतमर्त्थम्प्रद्योतयती'ति}{%
	\rYBh{1.01}{1.1}.  Cf.\ the reading {\dn सद्भूतमर्थम्} instead of {\dn स भूतमर्थम्} \nMaas\citep[5]{maas:2006}. See \pageline{BhT11s}--\pageline{BhT11e} in this edition.%
}%
\fubar ।  भूतार्त्थावगतिश्च \barfu विवेकख्यातिः।
%
त\vrtmm{स्माद्युक्त}{\Td\-M\-A\-\Me}{%
  \rdg{॰स्मात्द्युक्त॰}{\Tm}, 
  \rdg{॰\hlg{स्म}ाद्युक्तर॰}{L}}%
न्त\-दु\-पा\-य\-योगा\-नु\-शा\-स\-न\-मे\-वा\mms दौ\mme\ सूत्रितमिति॥

\llbl{p01_16}ननु च \barfu योगांगानुष्ठानाद्विवेक\varms{ख्यातिः। तत्र}{TM(em.)A\Me\vv{॰तिस्त॰}{TMA}, \vv{॰तिः, त॰}{\Me}, {\dn ख्यातिस्यत्र} L} वक्तव्यम`थ \varsm{योगांगानु}{\Tm LMA\Me, {\dn योगानु॰} \Td}शासनमि'ति॥

\llbl{p01_17}न। फलेनोपक्रमात्। योगांगानुष्ठानस्य हि फलं योग इति \barfu तदुपक्रमो \barfu युक्तः॥

\llbl{p01_18}\varsm{यद्येवं हानेनो}{\Tm LMA\Me, {\dn यद्येवमाह अनेनो॰} \Td}{\dn पक्रम्येत॥
  
\llbl{p01_19}न। तस्योपेयत्वात्। उपेयमेव हि तत्, योगः पुन स्वांगानामुपेयश्चोपायश्च विवेकख्यातेर\-%
\tstwll{animaadi}{अणिमादि॰}%
णिमादि\donote{animaadi}{%
	See \rYS{3.45}{3.45}: {\dn ततो \barfu ऽणिमादिप्रादुर्भावः \barfu कायसंपत्तद्धर्मानभिघातः॥} and the 
	\rnnYBh{3.45}{3.45} on it: {\dn तत्राणिमा \barfu भवत्यणुः। लघिमा लघुर्भवति।} \ldots\  %{\dn महिमा महान्भवति। प्राप्तिरङ्गुल्यग्रेणापि स्पृशति चन्द्रमसम्। प्राकाम्यमिच्छानभिघातः। भूमावुन्मज्जति निमज्जति यथोदके। वशित्वं भूतभौतिकेषु वशी भवत्यवश्यश्चान्येषाम्। ईशितृत्वं तेषां प्रभवाप्ययव्यूहानामीष्टे। यत्र कामावसायित्वं सत्यसंकल्पता यथा संकल्पस्तथा भूतप्रकृतीनामवस्थानम्।} %
	{\dn एतान्यष्टावैश्वर्याणि} \citep[164--5]{ys:anss}.%
}%
प्रा\mms प्ते\mme श्चे%
\varms{त्य\-भ्यर्हिततरः}{%
	\Tm\Tdac \Me(conj.?), %
	{\dn ॰त्यभ्यक्ततरः} \Tdpc, %
	{\dn ॰त्यह्यहिततरः} L, %
	{\dn ॰त्य \ccs इ\cce \ins हि\ine ततरः} M, %
	{\dn ॰त्यहिततरः} A, %
	{\dn ॰\eil त्यभ्य\eir र्हिततरः} \Me%
}%
, फल%
\vrtmm{द्वयनिमित्त}{conj.}{%
	\rdg{॰वन्निमित्त॰}{T}, %
	\rdg{॰(प/व)यनिमित्त॰}{L}, %
	\rdg{॰पयनिमित्त॰}{M ({\dn प} copiest marks as uncertain)}, %
	\rdg{॰पयनिमित्त॰}{A}, %
	\rdg{॰तदुपायरूप॰}{\Me}%
}%
त्वात्॥
  
\llbl{p01_20}उत्तर\varms{सूत्रार्त्थत्वाच्च}{LMA\Me\vv{॰र्थ॰}{MA\Me}, {\dn ॰सूत्रत्वाच्च} T}\fubar । न ह्य`थ योगांगानुशासनमि'त्युक्ते }\tst{`योगश्चित्तवृत्तिनिरोध' इति}{\rYS{1.02}{1.2}.} युक्तं \var{वक्तुम्। \barfu तदा}{M(em.)\vv{वक्तुं तदा}{M}A\Me, {\dn वक्तुन्ता\vv{॰क्तुम्। ता॰}{\Td}दा} TL} हि \tstsm{यमनियमादि सूत्र}{\rYS{2.29}{2.29}: {\dn यमनियमासनप्राणायामप्रत्याहारधारणाध्यानसमाधयो ऽष्टावङ्गानि।}.}मेव पठितव्यं स्यात्॥
  
\llbl{p01_21}अथ तदेव पठितव्यमिति \var{चेत््---न}{\Me(em.), {\dn चेत् न\vv{॰न्न}{MA} हि} TLMA}\fubar । \var{योगोपायाया}{TLMA, {\dn योगोपे\edl पा\edr याया} \Me} \barfu हानोपायभूताया \var{विवेकख्यातेः}{TMA\Me, {\dn विवेक\lost{2}तेः} L} फलस्य च हानस्य \varsm{सम्बन्ध}{T\-M\-A\-\Me, {\dn यम्बन्ध॰} L}\varmm{दर्शना\-}{TLMA, {\dn ॰\eil प्र\eir दर्शना॰} \Me}र्त्थ\varmm{त्वा`द्योग}{TL\Me(em.), {\dn त्वायोग॰} MA}श्चित्तवृत्तिनिरोध' इत्यस्य॥

\llbl{p01_22}\mms\varsm{कथम्पु}{\Td \Me(em.)\vv{\edl स\edr\  कथं पु॰}{\Me}, {\dn स कथम्पु॰} LMA}नरे\mme\varmm{तत्सम्बन्ध}{TMA, {\dn ॰तत्सम\-ब\-न्ध॰} L, {\dn ॰त\-त्सं\-ब\-न्ध\-प्र॰} \Me}दर्शना\vrtms{र्त्थम्}{TLA\Me\vv{र्त्थं}{TL}}{\rdg{॰र्थ॰}{M}}, \vrt{यावता}{LMA\Me}{\rdg{यावतां}{T}} निर्बीजसमाधिलक्षणाभिधानमेतत्॥
 
\llbl{p01_23}सत्यमेवम्। तथापि हानतदुपायसम्बन्ध\varms{म\-व\-द्यो\-त\-य\-दे\-व}{\Td(em.)\-A(em.)\-\Me(em.), {\dn ॰मपद्योतयदेव} \Tm, {\dn ॰मवद्योतयातिव} L, {\dn ॰मवद्योतयादेव} M} \varsm{सूत्र}{T\-L\-\Me(em.), {\dn सून॰} MA}न्निर्बीजसमा\mms धि\mds \vrtsm{लक्षणता}{LM\Apc \Me}{\rdg{लक्षण\ccs भिधानमेतत्सत्यमेवं तथापि\cce ता॰}{A}}मा\-द्य\-ते। न हि निरोध\-लक्षणात्स\-माधेरर्त्थान्तरं हानम्। किन्त्वे\mde ता\mme\varmm{वान्वि}{T, {\dn ॰वात्वि॰} L, {\dn ॰वता वि॰} MA\Me}शेषः---निरोधलक्षण\vrtms{समाधौ}{\Me(em.)}{\rdg{॰समाधौ न}{T}, \rdg{॰समाधेर्न्न\vv{॰र्न॰}{MA}}{L\-M\-A}, \rdg{॰समाधौ\edl धेर्न\edr }{\Me}} पुनःप्रवृत्तिः, हाने त्वात्यन्तिकी निवृ\mms त्ति\mds रिति। समाध्य\mte वस्थायान्तु हानाविशेष एव॥

%\begin{uncertain}% \com{ysbh13}{तथा चाह \ldots\ कैवल्यमि'ति}
\llbl{p01_24}तथा चाह \tstl{ys13}`\vrt{तदा}{\Tmpc \Td LMA\Me}{\rdg{त\ccs था\cce दा}{\Tm}} \var{द्रष्टुः}{\Tmac \Td LMA\Me, {\dn द्रष्टृ॰} \Tmpc} \vrt{स्वरूपे ऽवस्थानम्}{TLMA}{\rdg{स्वरूपेऽवस्थानम्' \eil इति\eir}{\Me}}'%
\donote{ys13}{\rYS{1.03}{1.3}}; %
\tstl{bh334part}`स्वरूपप्रतिष्ठा %
\vrt{वा}{em.}{\rdg{च}{$\Sigma$}} %
चितिशक्तिः' कैवल्यमिति%
\donote{bh334part}{Part of \rYS{4.34}{4.34} (the last sūtra): {\dn पुरुषार्थशून्यानां गुणानां प्रतिप्रसवः कैवल्यं स्वरूपप्रतिष्ठा वा चितिशक्तिरिति}.%  See the note on page~\pageref{ysbh13note}.
}\fubar । स्वरूप\vrtms{प्रतिष्ठत्वं हि}{\Me(em.)}{\rdg{॰प्रतिष्ठं हि}{TLMA}} \vrt{कैवल्यं हानम्}{conj.}{\rdg{कैवल्यमाह}{$\Sigma$}}\fubar ।%\end{uncertain}%\donote{ysbh13}{\come{ysbh13}}
\ ततश्च निर्बीजसमाधिना कैवल्यमेव शास्त्रार्त्थप्रत्य\mts यद्रढि\mde\-म्ने साक्षा\mme त्क्रियते॥
%\special{x:textcolor=000000}

\llbl{p01_25}तस्माद्यदुच्यते \tstwll{KAISC1}{कैश्चिन्निर्बीजसमाधिः कैवल्यसाधनमिति}\varsm{कैश्चिन्निर्बीज}{TL, {\dn कैश्चिद्बीज॰} MA, {\dn कैश्चित् \eil स\eir  बीज॰} \Me}समाधिः कैवल्यसाधनमिति\donote{KAISC1}{Cf.\ \Vacaspati\ on \rYS{1.03}{1.3}: {\dn निःश्रेयसस्य हेतुः समाधिरिति हि श्रुतिस्मृतीतिहासपुराणेषु प्रसिद्धम्॥} \citep[2,4--5]{ys:anss}.}, तन्न। किन्तर्हि ख्यातिरेव साधनमविद्यानिवृत्तिद्वारेण। अविद्यानिमित्तो हि \barfu बन्धः। तस्माद्यद्यपि योगानुशासनं सूत्रितम्, तथापि ख्यात्यर्त्थत्वाद्योगस्य तद्वारेण ख्यातिहानसम्बन्धमपि दर्शयतीति सुष्ठूच्यते। उत्तरसूत्रसम्बन्धार्त्थञ्च योगानुशासनमिति सूत्रयितव्यमिति॥

% \bigskip

\llbl{p01_26}\begin{uncertain}यद्यपि \mms\varsm{फलार्त्थि}{\Td\Lpc MA\Me, {\dn फलाति॰} \Lac}नान्तत्फल\mme साधनम्प्र\varmm{त्युत्प}{\Td LMA\Me, {\dn ॰त्युल्प॰} \Tm}न्नाकांक्षाणां साधनोपदेशो ऽप्यस्ति, \var{तथापि}{TMA\Me, {\dn \lost{1}थापि} L} न तत्प्रयोजनम्। सर्व्वस्य फलोद्देशेन \vrtsm{प्रवृत्ते}{\Tmpc \Td LMA\Me}{\rdg{प्रवृत्ति॰}{\Tmac}}स्तदेव प्रयोजनमिति॥\end{uncertain}

\begin{messwithp}{1000}{2pt}

%\newpage
\llbl{p01_27}\mll{ast1}%\com{anusasana}{योगानुशासनमिति \ldots शास्त्रम्}%
`योगानुशासनमि'ति---%
\fakepar
\tstwll{taiubhanu}{यथा शिष्यो ऽनुशिष्यते \ldots ॰नुशासनमित्युच्यते।}%
\begin{uncertain}यथा शिष्यो ऽनुशिष्यते विशिष्ट%
	\varmm{प्रवृत्तिनिवृत्ति\linelabel{nymdv}नियम}{%
		TMA\Me, 
		{\dn  ॰प्रवृत्तिनियम॰} L%
	}%
	द्वारेण \barfu तथा विशिष्टसाध्यसाधनतदंगनियममात्रसादृश्याद%
	\varmm{न्तेवास्य}{%
		\Td LA\Me, 
		{\dn ॰न्तेवास्या॰} \Tm, 
		{\dn ॰न्तेनास्य॰} M%
	}%
	नु\varmm{शासनवद्योगा}{%
		LMA\Me, 
		{\dn  ॰शासनविद्योगा॰} T%
	}%
	नुशासनमित्युच्यते।
\end{uncertain}%
%\resume{taiubhanu}
\donote{taiubhanu}%
{\labelh{TAIUBH111REF}Cf.\ \rTaiU{1.11.1}{1.11.1}: {\dn वेदमनूच्याचार्यो ऽन्ते\-वासि\-नम\-नुशास्ति} and the \rnnTaiUBh{1.11.1}{1.11.1} on it: {\dn वेदम\-नूच्ये\-त्येवमा\-दि\-कर्तव्यतो\-प\-देशारम्भः प्राग्ब्रह्म\-विज्ञाना\-न्नियमेन कर्त\-व्यानि \barfu श्रौत\-स्मार्त\-कर्माणी\-त्येवम\-र्थः। } {\dn \barfu अनुशासनश्रुतेः पुरुषसंस्कारार्थत्वात्। संस्कृतस्य हि विशुद्धसत्त्वस्यात्मज्ञानमञ्जसैवोत्पद्यते। `तपसा कल्मषं हन्ति विद्ययामृतमश्नुते'} (\rManu{12.104}{12.104}) {\dn इति हि \barfu स्मृतिः। \barfu वक्ष्यति \barfu च---`तप\-सा ब्रह्म विजिज्ञासस्व'} (\rTaiU{3.3.01}{3.3.1}) {\dn इति। \barfu अतो विद्यो\-त्पत्त्यर्थम\-नुष्ठेयानि कर्माणि। \barfu अनुशास्तीत्यनुशासनशब्दात्। अ\-नुशासना\-तिक्रमे हि \barfu दोषो\-त्पत्तिः। प्रागु\-पन्यासा\-च्च कर्मणाम्, केवल\-ब्रह्म\-विद्या\-रम्भाच्च। पूर्वं कर्माण्यु\-पन्यस्तानि। उदितायां च ब्रह्मविद्यायां, `अभयं प्रतिष्ठां विन्दते'} (\rTaiU{3.3.01}{3.3.1}) {\dn `न बिभेति कुतश्चन' `किमहं साधु नाकरवम्'} (\rTaiU{2.9.01}{2.9.1}) {\dn इत्येवमादिना कर्मनैष्किञ्चन्यं दर्शयिष्यतीत्यतो ऽवगम्यते---पूर्वोपचितदुरितक्षयद्वारेण विद्योत्पत्त्यर्थानि कर्माणीति। मन्त्र\-वर्णाच्च---`अविद्यया मृत्युं तीर्त्वा विद्ययामृतमश्नुते'} (\rIU{11}{11}) {\dn इति। ऋतादीनां \barfu पूर्वत्रोपदेशः \barfu आन\-र्थक्य\-परिहारा\-र्थः। इह तु \barfu ज्ञानो\-त्पत्त्यर्थ\-त्वा\-त्कर्तव्य\-नियमा\-र्थः। वेद\-मनूच्या\-ध्याप्या\-चार्यो ऽन्तेवासिनं शिष्यम\-नुशास्ति ग्रन्थ\-ग्रहणादनु पश्चाच्छास्ति तदर्थं \barfu ग्राहयती\-त्यर्थः। अतो ऽवगम्यते ऽधीत\-वेदस्य ध\-र्मजिज्ञासाम\-कृत्वा गुरु\-कुलान्न समा\-वर्तितव्यमिति। `बुद्ध्वा कर्माणि कुर्वीत'} (\Apastambhadharmasutra\ 9.21.5?) {\dn इति स्मृतेश्च।} \citep[420]{upa1}.%
} %
\barfu अ\-नुशिष्टिरनुशासनम्। \barfu यो\-गो ऽनुशिष्यते ऽने\varms{नास्मिन्निति}{\Td LM\Me, {\dn ॰नास्तिन्निति} \Tm, {\dn ॰नान्यस्मिन्निति} A} \barfu वा \mula{यो\-गा\-नुशासनं \var{शास्त्रम्}{\Td LMA\Me, {\dn सस्त्रम्} \Tm}}॥%\donote{anusasana}{\come{anusasana}}
\donote{ast1}{{\dn योगानुशासनं शास्त्रम्।}}

\end{messwithp}

\llbl{p01_28}\mll{athy1}%\com{athad}{अथेत्ययमधिकारा॰ \ldots\ प्रक्रियार्त्थो दृष्ट' इति}
\mula{अथेत्ययम\varmm{धिकारा}{LMA\Me, 
{\dn ॰धिकारो} T}र्त्थः}---अधिकार आरम्भः प्रस्तावः, अर्त्थो ऽभिधेयो ऽस्य। \tstl{sistasmrti}\pause{sistasmrti}<{\dn शिष्टस्मृतिप्रामाण्यात्}>\varsm{शिष्टस्मृति}{TL\Me, {\dn शिष्यस्मृति॰} MA}\varl{maanya}\pause{maanya}<॰माण्यात्॥ ननु>\-प्रामाण्यात्\donote{sistasmrti}{The \emph{śiṣṭasmṛti} here is the \VMBh.  See the beginning of the \VMBh: {\dn अथ शब्दा\-नुशासनम्। अथेत्ययं शब्दो ऽधिकारार्थः प्रयुज्यते।}.}॥\donote{athy1}{{\dn \barfu अथेत्ययमधिकारार्थः।}}\fakepar
%
\llbl{p01_29}ननु\resume{maanya}\donote{maanya}{TM(em.)A\Me\vv{॰ण्यान्ननु}{\Tm MA}, \rdg{॰प्रामण्यन्ननु}{L}} चानन्तर्य्यार्त्थस्याथशब्दस्य स्मरन्ति \barfu शिष्टाः। तथा \vrt{चाह}{\Tmpc \Td LMA\Me}{\rdg{चा\ccs न\cce\ins ह\ine}{\Tm}}---\tst{`वृत्तादनन्तरस्य प्रक्रियार्त्थो दृष्ट' \barfu इति}{\rSBh{1.1.1}. See also \rBSBh{1.1.01}{1.1.1}: {\dn अत्राथशब्द आनन्तर्यार्थः परिगृह्यते, नाधिकारार्थः, ब्रह्मजिज्ञासाया अनधिकारत्वात। मङ्गलस्य च वाक्यार्थे समन्वयाभावात्। अर्थान्तरप्रयुक्त एव ह्यथशब्दः श्रुत्या मङ्गलप्रयोजनो भवति। पूर्वप्रकृतापेक्षायाश्च फलत आनन्तर्याव्यतिरेकात्। सति चानन्तर्यार्थत्वे यथा धर्मजिज्ञासा पूर्ववृत्तं वेदाध्ययनं नियमेनापेक्षते, एवं ब्रह्मजिज्ञासापि यत्पूर्ववृत्तं नियमेनापेक्षते तद्वक्तव्यम्। स्वाध्यायनान्तर्यं तु समानम्। \ldots} \citep[47--51]{bsbh}.}\fubar ॥%\donote{athad}{\come{athad}}

\llbl{p01_30}न। सर्व्वत्रारम्भार्त्थत्वात्, आनन्तर्य्यस्य च गम्यमानत्वात्। \tstwll{PITAPUTRA}{यथा `पुत्र' इत्युक्ते \ldots\  पुत्रशब्दस्यार्त्थः \barfu पिता}%\com{putreti}{यथा `पुत्रे'त्युक्ते \ldots पुत्रस्यार्त्थः पिता}
यथा `पुत्र' \barfu इत्युक्ते \barfu पिता गम्यते, \var{न च}{\Tmpc \Td LMA\Me, {\dn न व} \Tmac} \var{पुत्रशब्दस्यार्त्थः}{LMA\Me\vv{\dn ॰र्थ्ः}{MA\Me}, {\dn पुत्रस्या}[\mist{1}]{\dn र्त्थः} \Tm, {\dn पु\-त्र\-स्या\-र्थः} \Td} \barfu \mms पि\mme ता।%\donote{putreti}{\come{putreti}}
\labelh{PTRPITR}\donote{PITAPUTRA}{Cf.\ \rBSBh{2.2.17}{2.2.17}: {\dn यथैको ऽपि सन्देवदत्तो लोके स्वरूपं संबन्धिरूपं \barfu चापेक्ष्यानेकशब्दप्रत्ययभाग्भवति---मनुष्यो \barfu ब्राह्मणः \barfu श्रोत्रियो \barfu वदान्यो \barfu बालो \barfu युवा \barfu स्थविरः \barfu पिता पुत्रः पौत्रो भ्राता जामातेति, यथा चैकापि सती रेखा स्थानान्यत्वेन निविशमानै\-कदश\-शत\-सहस्रा\-दि\-शब्द\-प्रत्यय\-भेदम\-नुभवति} \citep[521]{bsbh}.  This in turn alludes to \rYBh{3.13}{3.13}: {\dn यथैका रेखा शतस्थाने शतं दशस्थाने दशैका च एकस्थाने। यथैवैकत्वे ऽपि स्त्री माता चोच्यते दुहिता च स्वसा चेति।}.}
 \varsm{तथे}{\Td MA\Me, {\dn तथै॰} \Tm L}हा\mms प्या\mme रम्भ एवाभिधीयते, आनन्तर्य्यन्तु प्रतीयते। \var{अत एव च}{\Td\Lpc MA\Me, {\dn अत वेच} \Tm, {\dn अत एचव} \Lac} `वृत्तादनन्तरस्य प्रक्रियार्त्थो दृष्ट' इत्युच्यते। \barfu \varsm{यदि चा}{\Tmpc \Td\LMA\Me, {\dn यदि वा॰} \Tmac}न\-न्तर्य्या\-र्त्थो \barfu ऽभविष्यत्तदा `वृत्तादा\varms{नन्तर्य्यार्त्थे}{em., \rdg{॰नन्तर्य्ये र्त्थे}{\Tmac}, {\dn ॰नन्तर्यो र्त्थे} \Tmpc, {\dn ॰नन्तर्यो ऽर्थे} \Td, {\dn ॰नन्तर्य्यर्त्थे} L, {\dn ॰नन्तर्यो ऽर्थे} M, \rdg{॰नन्तर्ये ऽर्थे}{A\Me}} \var{प्रक्रियाया}{\Td\LMA\Me\vv{॰यायाः}{\Me}, {\dn प्रकियाया} \Tm}' इत्यवदिष्यत्॥
 
\llbl{p01_31}अनन्तरस्येति चानन्तरभाविन एवाभिधानम्।
अपि चैवं हि स्मृतिः---%
\labelh{VMBH1138_E}%
\tstl{mbh1138}\pause{mbh1138}<{\dn किञ्चिदव्ययं}\sel{\dn विद्यत इति}>%
`किञ्चिदव्ययं वि\-भ\-\mms \ktya \-\mme%
\varms{र्त्थप्रधानं}{%
	\Td\LMA\Me, 
	{\dn ॰र्त्थं प्रधानं} \Tm%
} %
किञ्चित्क्रियाप्रधानम्। उच्चै\-{}र्नीचैरि\-{}त्यादि विभ\ktya र्त्थ\-प्रधानम्। हिरु\-क्पृथ\-गि\-त्यादि %
\var{क्रियाप्रधानम्}{%
	T\LMA, 
	{\dn क्रियाप्रधानम् \eil इति\eir} \Me%
}%
\fubar । \fubar न \fubar चैत\-द्व्य\-ति\-रे\-के\-णा\-%
\varmm{व्य\-या\-ना}{%
	T\-M(em.)\-A\-\Me, 
	{\dn ॰प्रयाना॰} L%
}%
म%
\vrtms{र्त्थो}{\Td\LMA\Me}{%
	\rdg{॰र्त्थ॰}{\Tm}%
	} %
\var{विद्यत' इति}{em. {\dn वि\-द्य\-ते} $\Sigma$.}।\donote{mbh1138}{\rnnVMBh{1.1.38 or 1.2.47}{1.1.38 or 1.2.47} ad \rAAP{1.1.38}{1.1.38} or \rnAAP{1.2.47}{1.2.47}.  Cf.\ the readings in Kiel\-horn's third edition (p.\,95): {\dn किं\-चि\-दव्ययं विभ\ktya र्थप्रधानं किंचित्क्रियाप्रधानम्। उच्चैर्नीचैरिति वि\-भ\-\ktya\-र्थ\-प्र\-धा\-नं हिरु\-क्पृथगि\-ति क्रियाप्रधानम्। त\-द्धितश्चापि कश्चिद्विभ\ktya र्थप्रधानः \barfu कश्चित्क्रियाप्रधानः। तत्र यत्रेति विभ\ktya र्थप्रधानो नाना विनेति \barfu क्रियाप्रधानः। न चैतयोरर्थयोर्लिङ्गसंख्याभ्यां योगो ऽस्ति॥} and (p.~223) 
{\dn अव्ययं हि किंचिद्विभ\ktya र्थप्रधानं किंचित्क्रियाप्रधानम्। उच्चैर्नीचैरिति विभ\ktya र्थप्रधानं हिरुक्पृथगिति क्रियाप्रधानम्। तिबन्तं चापि किंचिद्विभ\ktya र्थप्रधानं किंचित्क्रियाप्रधानम्। काण्डे कुद्ये इति विभ\ktya र्थप्रधानं रमते ब्राह्मणकुलमिति क्रियाप्रधानम्। न चैतयोरर्थयोर्लिङ्गसंख्याभ्यां योगो ऽस्ति।}
} तत्रा\-थ\-शब्द\-स्या\-नन्तर्य्या\-र्त्थ\-त्वे सति \var{वि\-भक्ति\-श्रवणं युक्तम्}{T, {\dn वि\-भक्ति\-श्रय\-ण\-म\-युक्तं} L, {\dn वि\-भक्ति\-श्रवण\-म\-युक्तं} M, {\dn वि\-भक्ति\-श्रवण\-म\-प्यु\-क्तं} A, {\dn विभ\ktya श्रवणम्। विभक्तिश्रवणमयुक्तम्} \Me}, \var{सत्वप्रधानत्वात्}{\Tm L\Me\vv{॰सत्त्व॰}{\Me}, {\dn स त्वप्रधानत्वात्} \Td MA}\fubar । \barfu आरम्भक्रिया\varms{र्त्थत्वे तु सति}{em., {\dn ॰र्त्थत्वे तु स तु} \Tm L, {\dn ॰र्थत्वे तु} \Td, {\dn ॰र्थत्वे तु न} M\-A\-\Me} \vrtsm{वि\-भ\-\ktya}{\Tm\-L}{\rdg{विभ\ktya र्थ॰}{\Td MA}, \rdg{विभ\ktya\edl र्थ\edr ॰}{\Me}}\vrtmm{श्रवणं यु}{T}{\rdg{॰श्रवणमयु॰} L\-M\-A\-\Me\-\vv{॰श्रव\ccs न\cce णमयु॰}{L}}क्तम्, अस\mts त्वप्रधानत्वा\mte दिति॥

\llbl{p01_33}\labelh{TSMDTH}तस्मादथशब्दः प्रस्तावार्त्थ \vrtsm{एवेति युक्त}{TMA\Me}{\rdg{एवेति\lost{2}}{L}}म्॥

\llbl{p01_34}`अ\vrtmm{थेत्यधिका}{TMA\Me}{\rdg{॰[\hlg{थे}]त्यधिका॰}{L}}रार्त्थ' \barfu इति---स्वरूपा\varmm{वद्योतना}{TMA\Me, {\dn ॰वद्योतन\ccs ा\cce ॰} L}र्त्थ इति\mms \var{श\mme ब्दः}{\Td\LMA\Me, {\dn \lost{1}ब्दाः} \Tm}। % need checking (\Tm)!!!
यथा `गौरि'त्याहेति। प्र\vrtmm{सिद्धो ऽप्य\-थ}{TLM\Me\vv{॰द्धो प्य॰}{\Tm L}}{\rdg{॰सिद्ध्येप्यथ॰}{A}}\-शब्दार्त्थो ऽधिक्रियमाण\-\varms{विषये}{\Td\LMA\Me, {\dn विषय॰} \Tm} शिष्यबुद्धिसमाधानार्त्थं कीर्त्त्यते। \barfu सिद्धो \mts\vrt{ह्यय\-न्न्या\-यः \barfu श्रु\-तौ\mme}{\Me(em.)}{\rdg{ह्ययन्न्या\lost{2}तो}{L}, \rdg{ह्ययन्यायः \barfu श्रुतो}{M}, \rdg{\ccs ह्य\cce यं \barfu ना\ccs न्य\cce \ins य\ine ः। \barfu श्रुतो}{A}} % need checking!!!
\tstl{bau244}\pause{bau244}<व्याख्यास्यामि\sel निदिध्यायस्व>`\vrt{व्या\mde\-ख्या\-स्या\-मि}{\Td MA\Me}{\rdg{\lost{3}या\-ख्या\ccs\hlg{य्य}ाव्या\-दि\cce \ins स्या\ine मि}{\Tm}, \rdg{व्याख्यास्यास्यादि}{L}} ते \mds व्या\mms च\mte \vrtms{क्षाणस्य}{T\-\Me}{\rdg{॰क्ष\-ण\-स्य}{L\-M\-A}} तु मे \var{निदिध्यासस्वे\quotec ति}{\Tmpc\Me, {\dn निदिध्यासर्त्थेति} \Tmac, {\dn निदिध्यासस्येति} \Td, {\dn नि\-दि\-ध्या\-स\-ञ्चे\{॰सं चे॰ \textrm A\dn\}ति} \LMA}\donote{bau244}{\rBAU{2.4.04}{2.4.4} (\rnBAU{4.5.05}{4.5.5}).}\fubar ॥
  
\llbl{p01_35}\mll{khp1}कः \barfu पुनर्य्योगो यस्यानुशासनम्प्रस्तुतमित्याह---\linelabel{dpatha}\mula{योगः समाधिरि}ति। \barfu %
\tstwll{TWODP}{`यु\-जि\-{}र्य्यो\-गे' \ldots\ `युज समाधौ'}%
समाधिग्रहणा%
% \tstms{`द्यु\-जि\-{}र्य्यो\-ग' इति}{Dhātupāṭha 7.7.} निवर्त्त्य %
`द्यु\-जि\-{}र्य्यो\-ग' इति निवर्त्त्य %
%\tstl{dhp468}\pause{dhp468}<`युज समाधावि'त्य॰>%
`युज \var{समाधावि'त्ययं}{\Td MA\Me, {\dn समाधामित्ययं} \Tm L}%
%\donote{dhp468}{Dhātupāṭha 4.68.} %
\barfu गृहीतः।%
\donote{TWODP}{Dhātupāṭha 7.7 and 4.68.}
\linelabel{dpathaend} योगः समाधानम्॥\donote{khp1}{{\dn योगः \barfu समाधिः।}}

\llbl{p01_36}ननु च \mts `योगः स\mte माधिरि'\vrtms{त्ये\-व\-म\-न्ते\-न}{TL\Me}{\rdg{॰त्ये\-व\-म\-र्थे\-न}{\Mac}, \rdg{॰त्येव मतं तेन}{\Mpc A}} व्याख्याते सूत्रे `स च सार्व्वभौम' इत्यादि \varsm{भाष्यम\-सम्ब}{T\Mpc A\Me, {\dn भाष्यसम्ब॰} L, {\dn भाष्यं सम्ब॰} \Mac}न्ध\-न्दृश्यते? 

\begin{messwithp}{1000}{2pt}

\llbl{p01_37}नैष \barfu दोषः। इह `समाधिरि'त्युक्ते समाधीयमानापेक्षत्वात्; समाधीयमानानाञ्च बहुत्वात्; किमात्मा समाधीयते, किं \barfu वा
\varsm{शरीर}{\Td\LMA\Me, {\dn शारीर॰} \Tm}म्, \barfu आ\mts होस्विदिन्द्रियाणी\mte ति %\com{bahuvidaam}{बहुविदां विप्रतिपत्तेः}
\vrt{बहुविदां}{T\LMA}{\rdg{बहु\edl विदां\edr धा}{\Me}} \barfu विप्रतिपत्तेः%\donote{bahuvidaam}{\come{bahuvidaam}}
; \barfu स्व\-वि\-शे\-ष\-णाकांक्षत्वाच्च \barfu समाधेः॥

\end{messwithp}

\llbl{p01_38}\mll{ksyy1}कस्यायं किंविशेषणक इत्यनुषक्ते 
\vrt{प्रश्न}{em.}{\rdg{प्रश्ने}{\all}}
इदमाह---\mula{स च सार्व्वभौमश्चित्तस्य धर्म्म} इति। 
चित्तस्य धर्म्मो नात्मादीनाम्। तच्च चित्तं स्वयमेव समाधीयते 
\tstl{trentara}\pause{trentara}<{\dn कर्त्रन्तर॰}>\varsm{क\mms र्त्रन्त\mme र}{\Td, {\dn त्रेन्तर॰} L, {\dn \mist{5}न्तर॰} M, {\dn \mist{2}त्रेन्तर॰} A, {\dn \eil समाधात्र\eir न्तर॰} \Me}\donote{trentara}{Cf.\ note 34 of \citet{wezler:yvi1}\index{Wezler, Albrecht}.}निरपेक्षत्वात्॥\donote{ksyy1}{{\dn स च सार्व्वभौमश्चित्तस्य \barfu धर्म्मः।}}

\llbl{p01_39}\mll{bhmy1}काः पुनर्भूमय इत्यत आह---\mula{क्षिप्तम्मूढमि}त्यादयो \mula{भूमय} इति। 
\tst{क्षिप्तमिति कर्म्मकर्त्तरि \barfu निष्ठा। यथा भिन्नः कुसूलः स्वयमेवेति।}{%
  See \rVMBh{3.1.87, Vārttikas 5--7, esp.\ 7}{3.1.87, Vārttikas 5--7, esp.\ 7}: \ldots {\dn क्त। भिन्नः कुशूल इति कर्म। स  
  यदा स्वातन्त्र्येण विवक्षितस्तदास्य कर्मवद्भावः स्यात्। तस्य \barfu प्रतिषेधो वक्तव्यस्तस्मिन्प्रतिषिद्धे ऽकर्मकाणां भावे क्तो 
  भवतीति भावे क्तो यथा स्यात्। भिन्नं कुशूलेनेति। क्त॥} \citep[68]{mbh3}.%
} 
\mula{क्षिप्तम}निष्टविषया\vrtms{सञ्जनेन}{\Td\-L\-M\-A\-\Me}{\rdg{॰सञ्जनेने}{\Tm}} स्तिमितम्; 
\mula{मूढ}न्निर्व्वि\-{}वेकम्; \mula{विक्षिप्त}न्नानाक्षिप्तं कर्म्मकर्त्तर्य्येव; विवेकाक्षमञ्च वि\-क्षिप्त\-त्वादेव; 
\mula{एकाग्र}न्तुल्य\-प्र\mms त्य\mme\-य\-\vrtms{प्रवाहम्}{\Me(em.)}{\rdg{॰वाहम्}{T\LMA}}; \mula{निरुद्धमिति} प्रत्ययशून्यम्---चित्तम्॥\donote{bhmy1}{{\dn क्षिप्तम्मूढं विक्षिप्तमेकाग्रन्निरुद्धमिति \barfu भूमयः।}}

\llbl{p01_40}ननु च भूमि\mms षु\mme\ धर्म्मेषु विवक्षितेषु किमर्त्थं क्षिप्तमित्यादिना धर्म्म्युच्यते?

\llbl{p01_41}नैष \barfu दोषः। धर्म्मिणा धर्म्म एवोपदिश्यते। 
धर्म्माणान्धर्म्मिविषय\varms{त्वात्। \barfu यथा}{\Td\LMA\-\Me\vv{\dn त्वाद्यथा}{\Td\LMA}, {\dn त्वात् \barfu द्यथा} \Tm} गोत्वे किं लिंगमिति पृष्टे विषाणी ककुद्मान्प्रान्तवालधिरिति धर्म्मिणा धर्म्म 
एवोपदिश्यते। तस्माद्विक्षेपादयश्चित्तस्य \var{भूमयो}{\Tmpc\Td\Lpc MA\Me, {\dn भूमय} \Tmac, {\dn भूयो} \Lac} धर्म्मा \barfu इत्यर्त्थः॥ 
  
\llbl{p01_42}ननु च भूमीनाञ्चित्तधर्म्मत्वे समाधेश्च, कथम्भूमिभिः समाधिराधारत्वेन विशेष्यते 
\var{सार्व्वभौम}{\Me(em.),  {\dn सार्व्वभूम} T\LMA} \barfu इति---\var{न}{TLM, om.\ A\Me}\fubar । 
\barfu सामान्य\-भूत\-त्वात्समाधेः, भूमीनाञ्च विशेषत्वात्---\barfu यथा 
\vrtsm{क्षिप्त}{\Me(em.)}{\rdg{विक्षिप्त॰}{\Td MA}, \rdg{विक्षिप्ता॰}{\Tm L}}न्तिष्ठति, मूढन्तिष्ठति, \vrtsm{विक्षिप्त}{\Me(em.)}{\rdg{क्षिप्त॰}{T\LMA}}मेकाग्र\vrtmm{न्निरुद्ध}{conj.}{\rdg{\om}{$\Sigma$}}ञ्चेति क्षिप्तादिभूमिषु 
स्थितिरनुवर्त्तते सामान्यम्। %
\linelabel{smdhsth}%
\barfu स्थितिश्च \barfu समाधिः। तस्मात्सामान्यरूपः सर्व्वासु भूमि\mms षु\mme\ 
\vrt{साधारण्येन}{em.}{\rdg{साधान्येन}{\Tmac}, \rdg{प्राधान्येन}{\Tmpc \Td\LMA\Me}} वर्त्तत इति॥

%\comml{aa5141}
\llbl{p01_43}\tst{`सर्व्वभूमिपृथिवी'त्यैश्वर्य्यार्त्थीयो \barfu ऽण्प्रत्ययः। `अनुशतिकादि'त्वादुभयपदवृद्धिः}{See \rAAP{5.1.41--2}{5.1.41--2}: {\dn सर्वभूमिपृथिवीभ्यामणञौ। \barfu तस्ये\-श्व\-रः।} and  \rAAP{7.3.20}{7.3.20}: {\dn अनुशतिकादीनां च}.}%\donote{aa5141}{\come{aa5141}}
॥ 
% the note should be more elaborate!!!

\llbl{p01_44}\tstwll{bhuumis}{अन्ये पुनर्बाह्याध्यात्मिकान्संयमविषयान्भूमय
इत्याचक्षते।}अन्ये\labelh{ANYEPUNARedn}%\com{samayama}{अन्ये \ldots संयमो भवति}
\ पुनर्बाह्याध्यात्मिकान् 
\vrtsm{संयम}{T\Me}{\rdg{समयम॰}{\LMA}}% CHECK!
विषयान्भूमय इत्याचक्षते।\donote{bhuumis}{Cf.\ \rYS{3.06}{3.6}: {\dn तस्य भूमिषु \barfu विनियोगः} and the \rnnYBh{3.06}{3.6} on it. The objects (\visaya) of \samyama\ are described in \rYS{3.16}{3.16}\,ff.}
तेषां \tst{`विक्षिप्ते चेतसी'ति}{See\ p.~\pageref{VKSPTCTS1}} सामानाधिकरण्यानुपपत्तिः, स्ववचनविरोधश्च। 
कथं? क्षिप्यते ऽस्मिन्निति क्षिप्तम्, मुह्यते ऽस्मिन्निति मूढम्, \vrt{विक्षिप्यते ऽस्मिन्निति \barfu च विक्षिप्तम्}{conj.}{\rdg{\om}{$\Sigma$}}\fubar । यदा चैवन्न तदा \barfu \vrtsm{संयम}{\Tmpc\Td\Me}{\rdg{समयम॰}{\Tmac LMA}}\-%
विषयता, तत्र क्षेपाद्यसम्भवात्। एकाग्रावस्थायां हि \vrt{संयमो}{T\Me}{\rdg{समयमो}{LMA}} भवति॥%\donote{samayama}{\come{samayama}}

\llbl{p01_45}\vrt{किञ्चान्यत्}{T\Me(em.)}{{\dn किञ्चान्य} LMA}। \tst{क्षिपेश्च ध्रौव्याद्यर्त्थाभावादधिकरणे \mts निष्ठा\vrt{प्रत्यया\mte भावः}{LMA\Me}{{\dn \lost{3}भावः} \Tm, {\dn \mist{3}मा\-नः} \Td}}{See \rAAP{3.4.76}{3.4.76} {\dn क्तो ऽधिकरणे च \barfu ध्रौव्य\-गति\-प्रत्य\-व\-साना\-र्थेभ्यः॥}}॥

\llbl{p01_46}किञ्च निरुद्धे चेतसि \vrtsm{संयम}{T\Me(em.)}{\rdg{समयम॰}{LMA}}स्यापि \barfu तत्राभावाद्वृत्तिशून्य\vrtmm{त्वात्संयम}{T\Me(em.)}{\rdg{॰त्वात्समयम॰}{LMA}}विषयाभावः। न हि निरोधे विषयविशेषे चित्तन्निरुध्यते। तदा हि विषयित्वान्निवर्त्तते। परिगणनानुपपत्तिश्च। न हि \vrt{क्षि\-प्ता\-दयः}{\Me(em.)}{\rdg{विक्षिप्तादयः}{TLMA}} पञ्चैव \vrt{संयमस्य}{T\Me(em.)}{\rdg{समयस्य}{LMA}} विषयास्तेषामानन्त्यात्॥

\llbl{p01_47}ननु \vrt{चान्यथा}{T\-M\-A\-\Me}{\rdg{चान्यधा}{L}} योगं केचिदिच्छन्ति? तथा चाहुः---
\begin{quote}\firmlist
%\com{vs}{इन्द्रियमनो॰ \ldots\ संयोगः}%
\llbl{p01_48}%
\tstwll{VSFIVETWO}{इन्द्रियमनोर्त्थसन्निकर्षात्सुखदुःखे। तदनारम्भ आत्मस्थे मनसि सशरीरस्य सुखदुःखाभावः प्राणमनोविनिग्रहापेक्षः \barfu संयोगः}%
\str{इन्द्रियमनोर्त्थसन्निकर्षात्सुखदुःखे। तदनारम्भ आत्मस्थे मनसि %
\vrt{स\-श\-री\-र\mms स्य\mme}{\Td\Me(em.)}{%
	\rdg{स \barfu सशरीर[\mist{1}]}{\Tm}, 
	\rdg{शरीरस्य}{LMA}, 
	\rdg{\eil स\eir शरीरस्य}{\Me}%
} 
सुखदुःखा\-%
\vrtms{भावः}{TL\Me}{%
	\rdg{भावं}{MA}%
} %
\barfu प्राण\-मनो\-वि\-निग्रहा\-पेक्षः \vrt{संयोगः}{TLMA}{%
	\rdg{संयोगो \eil योग\eir\  इति}{\Me}%
}%
॥%
}\donote{VSFIVETWO}{\rVS{5.2.16--7}.  See \citet{wezler:yviandvs}\index{Wezler, Albrecht}.  Cf.\ \Carakasamhita\ \citep{caraka} Śarīra\-sthā\-na 2.138--139:
{\dn
    आत्मे\-न्द्रिय\-मनोऽर्थानां सन्नि\-कर्षा\-त्प्रवर्तते।
    सुख\-दुःख\-मनारम्भादा\-त्म\-स्थे मनसि स्थिरे॥
    नि\-वर्तते तदु\-भयं वशि\-त्वं चो\-पजायते।
    स\-शरीर\-स्य योग\-ज्ञाः तं \barfu योग\-मृषयो \barfu विदुः॥}; % \end{verse}
\rNBh{4.2.38}: {\dn स तु प्रत्याहृत\-स्ये\-न्द्रिये\-भ्यो मन\-सो धारके\-ण प्र\-यत्नेन धार्य\-माण\-स्या\-त्मना \barfu संयोग\-स्तत्त्व\-बुभुत्सा\-विशिष्टः। सति हि तस्मि\-न्निन्द्रिया\-र्थेषु बुद्ध\-यो नो\-त्पद्यन्ते, तद\-भ्यास\-वशा\-त्तत्त्व\-बुद्धि\-रुत्पद्य\-ते॥}; \PDhS\ \citep[187]{pdhs}: {\dn अस्म\-द्विशिष्टानान्तु योगि\-नां युक्तानां योगज\-धर्मा\-नु\-गृही\-तेन मनसा स्वात्मा\-न्तरा\-काश\-दिक्काल\-परमाणु\-वायु\-मनस्सु तत्समवेत\-गुण\-कर्म\-सामान्य\-विशेषेषु सम\-वाये चा\-वितथं स्वरूप\-दर्शन\-मुत्पद्यते। वियुक्ता\-नाम्पुन\-श्चतुष्टय\-सन्निकर्षा\-द्योगज\-धर्मा\-नुग्रह\-सा\-मर्थ्यात्सूक्ष्म\-व्यवहित\-विप्रकृष्टेषु प्रत्य\-क्षमुत्पद्यते। तत्र सामान्य\-विशेषेषु स्वरूपा\-लोचन\-मात्र\-म्प्रत्यक्ष\-म्प्रमाण\-म्प्रमेया द्रव्यादयः पदार्थाः प्रमाना\-त्मा प्रमिति\-र्द्रव्यादि\-विषय\-ञ्ज्ञानम्। सामान्य\-विशेष\-ज्ञानो\-त्पत्तावविभक्त\-मालोचन\-मात्र\-म्प्रत्यक्ष\-म्प्रमाण\-मस्मिन्नान्यत्प्रमाणा\-न्तरम\-स्त्य\-फलरूपत्वात्। अथ वा सर्वेषु पदार्थेषु चतुष्टय\-सन्निकर्षा\-दवि\-तथम\-व्यपदेश्यं यज्ज्ञान\-मुत्पद्यते तत्प्रत्यक्ष\-म्प्रमाणम्प्रमेया द्रव्यादयः पदार्थाः प्रमाता\-त्मा प्रमिति\-र्गुण\-दोष\-माध्यस्थ्य\-दर्शन\-मिति॥}; \PaSu\ \citep[6]{pasupatasutra}: {\dn किं तदि\-त्यु\-च्यते योगम्। \barfu अत्रा\-त्मेश्वर\-संयोगो \barfu योगः।}}%\donote{vs}{\come{vs}}

\llbl{p01_49}\linelabel{vscomstart}\mula{तदि}ति प्रकृतापेक्षम्। \linelabel{yosau}यो \barfu ऽसौ \mula{सुखदुःख}योरात्मे\mula{न्द्रियमनोर्त्थसन्निकर्षो} हेतुः, \barfu \linelabel{tasya}%
त\-स्या\-\mula{नारम्भ}स्तदनुत्पत्तिः। स कथम्भवति? \mula{आत्मस्थे मनसि} %
\vrt{नेन्द्रियस्थे}{TMA\Me}{%
    \rdg{ने[\mist{4}]}{L}}%
। \mula{स\-श\-री\-र\-स्या}\-विशीर्ण्णशरीरस्य। \barfu तदा %
\tst{`कारणाभावात्कार्याभाव' इति}{\rVS{1.2.1}.} %
सन्निकर्षा\-{}भावे \barfu \mula{सु\-ख\-दुः\-ख}\-%
\vrtmm{योरप्य\mula{भाव}}{T\Lpc MA\Me}{%
    \rdg{॰योरप्यवा॰}{\Lac}%
}ः। \vrt{तस्याम\-वस्था\-यां}{MA\Me}{\rdg{तस्यावस्थायां}{\Tm}, \rdg{तद\-व\-स्था\-यां}{\Td}, \rdg{तस्यामस्थायां}{L}} \barfu यो \barfu ऽसौ \barfu विभोरात्मनो \barfu मनसा \vrt{\mula{संयोगः}}{TMA\Me}{\rdg{सायोगः}{L}}, \linelabel{vscomsa}स \mula{प्राण\-मनो\-विनिग्रहा\vrtms{पेक्षः}{MA\Me}{\rdg{॰पेक्ष॰}{TL}}} \barfu संयोगविशेषो \barfu योगः\linelabel{vscomend}
\end{quote}

%\newpage
---\varl{ntoc1}\pause{ntoc1}<{\dn इति॥ अत्रो॰}>{इति॥%
\fakepar
\llbl{p01_50}अत्रो}\donote{ntoc1}{LMA\Me\vv{इत्यत्रो॰}{LMA}, \rdg{इत्यन्तो॰}{T}}च्यते---


\llbl{p01_51}\linelabel{aatmasthe}`आत्मस्थे मनसी'त्ययुक्तम्, \barfu \vrtsm{सर्व्वदात्म}{TLMA}{\rdg{सर्वदा आत्म॰}{\Me}}स्थत्वा\mms न्मन\mme सः॥

\llbl{p01_52}इन्द्रियादिसन्निकर्षा\vrtmm{नारम्भा}{\Tmpc\Td\Lpc MA\Me}{\rdg{॰दारम्भा॰}{\Tmac}, \rdg{॰\ccs द\cce नारम्भा॰}{L}}पे\vrtmm{क्षया`त्म}{LMA}{\rdg{॰क्षया आत्म॰}{T\Me}}\labelh{TEcorr1}\vrtms{स्थ' इति---इति}{em.}{\rdg{॰स्थ इति}{TLMA}, \rdg{॰स्थे इति}{\Me}} चेत्---`तदनारम्भ' इत्येव सिद्ध\-{}त्वा\-{}दा`त्मस्थे मनसी'त्यनर्त्थकं स्यात्॥

\llbl{p01_53}कि\mts ञ्चान्यत्---मु\mte क्तस्यापि मनस आत्मस्थत्वादिन्द्रियसन्निकर्षाभावाच्च योग\mms ः\mme\ प्राप्नोति।
\tstwll{atmpm1}{तस्यापि हि \barfu सर्व्वगतत्वात्}तस्यापि हि \barfu सर्व्वगतत्वा\donote{atmpm1}{See \PDhS\ \citep[70]{pdhs}: {\dn तथा चात्मेतिवचना\-त्परम\-मह\-त्परि\-माणम्}.  Cf.\ \rVS{7.1.28--9}: {\dn \barfu विभवान्महानाकाशः। \barfu तथाचात्मा।}.}\tstwll{mnsnt1}{मनसश्च नित्यत्वात्}न्मनसश्च नित्यत्वा\donote{mnsnt1}{See \rVS{3.2.2}: {\dn द्रव्यत्वनित्यत्वे वायुना व्याख्याते।}; \PDhS\ \citep[130]{pdhs}: {\dn नित्यं परमाणुमनस्सु तत्पारिमाण्डलुयम्}.}त्स\vrtmm{र्व्वदात्म}{LMA\Me}{\rdg{॰र्व्वन्त\vv{॰र्वं त}{\Td}दात्म॰}{T}}स्थत्वमेव॥ % VS 3.2.3 manas' nityatva; PDhS 70 tathaa caatmetivacanaatparamamahatparimaa.nam (VS 7.1.29 tathaacaatmaa)

\llbl{p01_54}\tstwll{ATMPRDS1}{किञ्चात्मनः \ldots\ मिथ्यात्वात्}किञ्चात्मनः प्रदेशाभावादात्मस्थ \vrt{इत्ययुक्तम्}{TMA\Me}{\rdg{इत्य[\mist3]}{L}}। %\ldots\ % ref?
न \barfu चा\mts प्युपचरितात्म\mte प्रदेशसंयोगः परमार्त्थस्य योगस्य हेतु \vrtsm{स्यात्। उप}{\Tm\Lpc MA\Me\vv{स्यादुप॰}{\Tm MA}}{\rdg{स्यात् अप॰}{\Td}, \rdg{॰त्वादुप॰}{\Lac}}\-च\-रित\-स्य मिथ्यात्वात्॥\donote{ATMPRDS1}{\labelh{BSBHTHREETHREE}Cf.\ \rnnNS{2.2.13–17}{2.2.13–17}/\rnnNBh{2.2.13–17}\ 2.2.13–17: {\dn अनित्यः शब्द इत्युत्तरम्। कथम्?---{\dnb आदिमत्त्वादैन्द्रियकत्वात्कृतकवदुपचाराच्च॥} \ldots\ इतश्च शब्द उत्पद्यते नाभिव्यज्यते---कृतकवदुपचारात्। तीव्रं मन्दमिति कृतकमुपचर्यते, तीव्रं सुखं मन्दं सुखं तीव्रं दुःखं मन्दं दुःखमिति, उपचर्यते च तीव्रः शब्दो मन्दः शब्द इति। \ldots\ 
% p.603
%vyañjakasya tathābhāvād grahaṇasya tīvramandatā rūpavaditi cenna abhibhavopapatteḥ---saṃyogasya vyañjakasya tīvramandatayā śabdagrahaṇasya tīvramandatā bhavati na tu śabdo bhidyate yathā prakāśasya tīvramandatayā rūpagrahaṇasyeti, tac ca na, evam abhibhavopapatteḥ---tīvro bherīśabdo mandaṃ tantrīśabdam abhibhavati na mandaḥ/ na ca śabdagrahaṇam abhibhāvakam, śabdaś ca na bhidyate, śabde tu bhidyamāne yukto 'bhibhavaḥ/ 
तस्मादुत्पद्यते शब्दो नाभिव्यज्यत इति। \ldots\ 
%
% p.604
% abhibhavānupapattiś ca vyañjakasamānadeśasyābhivyaktau prāptyabhāvāt---vyañjakena samānadeśe 'bhivyajyate śabda ity etasmin pakṣe nopapadyate 'bhibhavaḥ/ na hi bherīśabdena tantrīsvanaḥ prāpta iti/
%
% p.605
%aprāpte 'bhibhava iti cet śabdamātrābhibhavaprasaṅgaḥ---atha manyeta asatyām prāptāv abhibhavo bhavatīti, evaṃ sati yathā bherīśabdaḥ kañcit tantrīsvanam abhibhavati evam antikasthopādānam iva davīyaḥstho pādānān api tantrīsvanān abhibhaved aprāpter aviśeṣāt/ tatra kvacid eva bheryāṃ praṇāditāyāṃ sarvalokeṣu samānakālās tantrīsvanā na śrūyeran iti/ nānābhūteṣu śabdasantāneṣu satsu śrotrapratyāsattibhāvena kasyacic chabdasya tīvreṇa mandasyābhibhavo yukta iti/
%
% mandasyābhibhavo yukta iti/] p.606
%kaḥ punar ayam abhibhavo nāma? grāhyasamānajātīyagrahaṇakṛtam agrahaṇam abhibhavaḥ; yatholkāprakāśasya grahaṇārhasyādityaprakāśeneti
॥१३॥
%
%
%_____________________________________________________________________
%
%
%     ********************** NySBh_2,2.14 **********************
%
%
% p.607
%
%     NyS_2,2.14: 
{\dnb न, घटाभावसामान्यनित्यत्वान्नित्येष्वप्यनित्यवदुपचाराच्च॥१४॥}
%
\ldots\ 
%na khalu ādimattvād anityaḥ śabdaḥ/ kasmāt? vyabhicārāt/ ādimataḥ khalu ghaṭābhāvasya dṛṣṭaṃ nityatvam/ katham ādimān? kāraṇavibhāgebhyo hi ghaṭo na bhavati/ katham asya nityatvam? yo 'sau kāraṇavibhāgebhyo na bhavati na tasyābhāvo bhāvena kadācin nivartyata iti/ yad apy aindriyakatvāt, tad api vyabhicarati, aindriyakaṃ ca sāmānyaṃ nityaṃ ceti/ 
यदपि कृतकवदुपचारादिति, एतदपि व्यभिचरति; नित्येष्वनित्यवदुपचारो दृष्टः---तथा हि भवति वृक्षस्य प्रदेशः कम्बलस्य प्रदेशः, एवमाकाशस्य प्रदेशः आत्मनः प्रदेश इति भवतीति॥१४॥
\ldots\ 
%
%_____________________________________________________________________
%
%
%     ********************** NySBh_2,2.15 **********************
%
%
%% p.608
%
%     NyS_2,2.15: tattvabhāktayor nānātvasya vibhāgād avyabhicāraḥ ||
%
%nityam ity atra kiṃ tāvat tattvam? arthāntarasyānutpattidharmakasyātmahānānupapattir nityatvam, tac cābhāve nopapadyate/ bhāktaṃ tu bhavati yat tatrātmānam ahāsīd yad bhūtvā na bhavati na jātu tat punar bhavati, tatra nitya iva nityo ghaṭābhāva ity ayaṃ padārtha iti/ tatra yathājātīyakaḥ śabdo na tathājātīyakaṃ kāryaṃ kiṃcin nityaṃ dṛśyata ity avyabhicāraḥ//
%
%
%_____________________________________________________________________
%
%
%     ********************** NySBh_2,2.16 **********************
%
%
%% p.609
%yad api sāmānyanityatvād itīndriyapratyāsattigrāhyam aindriyakam iti---
%
%     NyS_2,2.16: santānānumānaviśeṣaṇāt ||
%
%nityeṣv avyabhicāra iti prakṛtam/ nendriyagrahaṇasāmarthyāc chabdasyānityatvam/ kiṃ tarhi indriyapratyāsattigrāhyatvāt santānānumānaṃ tenānityatvam iti //
%
%
%_____________________________________________________________________
%
%
%     ********************** NySBh_2,2.17 **********************
%
यदपि नित्येष्वप्यनित्यवदुपचारादिति। न---%
%
%     NyS_2,2.17: 
{\dnb कारणद्रव्यस्य प्रदेशशब्देनाभिधानात्॥१७॥}
%
नित्येष्वप्यव्यभिचार इति। एवमाकाशप्रदेश आत्मप्रदेश इति नात्राकाशात्मनोः कारणद्रव्यमभिधीयते यथा कृतकस्य। कथं ह्यविद्यमानमभिधीयते? \barfu अविद्यमानता च \barfu प्रमाणतो \barfu ऽनुपलब्धेः। किंतर्हि तत्राभिधीयते संयोगस्याव्याप्यवृत्तित्वम्। \ldots\ %
%परिच्छिन्नेन द्रव्येनाकाशस्य संयोगो नाकाशं व्याप्नोति, अव्याप्य वर्तत इति, तदस्य कृतकेन द्रव्येण सामान्यम्। न \barfu ह्यामलकयोः संयोग आश्रयं व्याप्नोति। %
\barfu सामान्यकृता च भक्तिराकाशस्य प्रदेश इति। \barfu अनेनात्मप्रदेशो \barfu व्याख्यातः।} \ldots\ \citep[594–613]{ns}; \SplitNote%
%\rBSBh{2.2.17}{2.2.17}: {\dn तथाण्वात्ममनसामप्रदेशत्वान्न \barfu संयोगः संभवति। \barfu प्रदेशवतो द्रव्यस्य \barfu प्रदेशवता द्रव्यान्तरेण संयोगदर्शनात्। कल्पिताः प्रदेशा अण्वात्ममनसां भविष्यन्तीति चेत्, न, अविद्यमानार्थकल्पनायां सर्वार्थसिद्धिप्रसंगात्।} \ldots\ \citep[521]{bsbh}; 2.3.53: {\dn अथोच्येत---विभुत्वे ऽप्यात्मनः शरीरप्रतिष्ठेन मनसा संयोगः शरीरावच्छिन्न एव आत्मप्रदेशे भविष्यति, अतः \barfu प्रदेशकृता व्यवस्थाभिसंध्यादीनामदृष्टस्य सुखदुःखयोश्च भविष्यतीति। तदपि नोपपद्यते। कस्मात्? अन्तर्भावात्। विभुत्वाविशेषाद्धि सर्व एवात्मानः सर्वशरीरेष्वन्तर्भवन्ति। तत्र न वैशेषिकैः शरीरावच्छिन्नो \barfu ऽप्यात्मनः प्रदेशः कल्पयितुं \barfu शक्यः। कल्प्यमानो ऽप्ययं निष्प्रदेशस्यात्मनः प्रदेशः \barfu काल्पनिकत्वादेव न पारमार्थिकं कार्यं नियन्तुं शक्नोति। \ldots} \citep[627]{bsbh}.%
}

\llbl{p01_55}\linelabel{indriyamano}किञ्चान्यत्। मनोर्त्थसन्नि\vrtmm{कर्षादि}{em.}{\rdg{॰कर्षाभाव इ॰}{$\Sigma$}}त्येव \vrtsm{सिद्धेरि}{\Me(em.)}{\rdg{सिद्धमि॰}{TLMA}, \rdg{सिद्धेः\edl द्धं\edr\ इ॰}{\Me}}न्द्रियग्रहणमनर्त्थकं \vrtsm{स्या}{TLA\Me}{\rdg{॰त्वा॰}{M}}त्॥\FootnotetextC{}{Cf.\ also \rBSBh{2.2.17}{2.2.17}: {\dn तथाण्वात्ममनसामप्रदेशत्वान्न \barfu संयोगः संभवति। \barfu प्रदेशवतो द्रव्यस्य \barfu प्रदेशवता द्रव्यान्तरेण संयोगदर्शनात्। कल्पिताः प्रदेशा अण्वात्ममनसां भविष्यन्तीति चेत्, न, अविद्यमानार्थकल्पनायां सर्वार्थसिद्धिप्रसंगात्।} \ldots\ \citep[521]{bsbh}; 2.3.53: {\dn अथोच्येत---विभुत्वे ऽप्यात्मनः शरीरप्रतिष्ठेन मनसा संयोगः शरीरावच्छिन्न एव आत्मप्रदेशे भविष्यति, \barfu अतः \barfu प्रदेशकृता व्यवस्थाभिसंध्यादीनामदृष्टस्य सुखदुःखयोश्च भविष्यतीति। तदपि नोपपद्यते। कस्मात्? अन्तर्भावात्। विभुत्वाविशेषाद्धि सर्व एवात्मानः सर्वशरीरेष्वन्तर्भवन्ति। तत्र न \barfu वैशेषिकैः \barfu शरीरावच्छिन्नो \barfu ऽप्यात्मनः प्रदेशः कल्पयितुं \barfu शक्यः। कल्प्यमानो ऽप्ययं निष्प्रदेशस्यात्मनः \barfu प्रदेशः \barfu काल्पनिकत्वादेव न पारमार्थिकं कार्यं नियन्तुं शक्नोति। \ldots} (ibid.:\ 627)%
%\citep[627]{bsbh}%
.}

\llbl{p01_56}अथेन्द्रियैर्व्विना मनोर्त्थसन्निकर्षा\mts भावात्तदना\mte रम्भ इति प्राप्त्य\vrtms{भावे}{\Td\LMA\Me}{\rdg{॰भावो}{\Tm}} प्रतिषेधो न युक्त इति चेत्---नैतदेवम्। विनेन्द्रियैर्म्मनोर्त्थसन्निकर्षो नास्तीत्यर्त्थात्सन्निकर्षप्रतिषेधं कुर्व्वता विषयप्राप्तिरिन्द्रियद्वारा \barfu प्रतिपादिता %
%\rule{0pt}{10pt} %
\mms 
स्यात्॥\linelabel{indriyamanoend}

\llbl{p01_57}\linelabel{sukhaduhkhe}\vrtsm{किञ्चा\mds न्य}{LMA\Me}{\rdg{किंत्वा\mist{1}}{\Td}}त्---सर्व्वप्राण\mde भा\mme जामिन्द्रियमनोर्त्थसन्निकर्षा\vrtmm{त्सुखदुःख}{\Tmpc \Td\LMA\Me}{\rdg{॰त्सुख॰}{\Tmac}}प्रतिलम्भस्य \vrtms{प्रसिद्ध\-त्वा}{\Td\LMA\-\Me}{\rdg{प्रसिद्धत्वप्रसिद्धत्वा॰}{\Tm}}\-त्, % ref. 
 आ\-\vrtmm{त्ममनः}{em.}{\rdg{॰त्मनः॰}{$\Sigma$}}सम्बन्धस्य च नि\-त्य\-त्वात्, सुखदुःखाभावो योग इत्येव \vrt{सिद्धेः \barfu सूत्रशेषो}{\conj}{\rdg{\ccs सि सू\cce सिद्धे सूत्रविशेषो }{\Tm}, \rdg{सिद्धे सूत्रविशेषो}{ \Td}, \rdg{सिद्धेस्तत्र \ilg षो}{L}, \rdg{सिद्धेस्तत्र शेषो }{MA}, \rdg{सिद्धेः, तत्र \eil सूत्र\eir शेषो}{\Me}} \barfu ऽनर्त्थक स्या\-\mms त्॥

\llbl{p01_58}अथापि स्यादि\mms न्द्रि\mme य\mms मनो\mds र्त्थसन्नि\vrtmm{कर्षग्रहणा}{L\Mpc A\Me}{\rdg{॰कर्ष\il ग्रहणा\ir ॰}{M}}दृते\mde\  सुखदुःखाभाव\mme\ इत्येतावति सति मुक्तस्या\-{}पि \barfu सुखदुःखाभावाद्यो\vrtmm{गित्वप्र}{\Me(em.)}{\rdg{॰ग\ccs स्य\cce त्वप्र॰}{\Tm}, \rdg{॰गस्य त्वप्र॰}{\Td}, \rdg{॰गत्वप्र॰}{LMA}}संग इति तन्निवृत्यर्त्थं सूत्र\vrtmm{शेषग्रहण}{\Td\LMA\Me}{\rdg{॰शेषप्रग्रहण॰}{\Tm}}मिति---

\llbl{p01_59}तच्च न, मुक्तस्य सुखदुःखप्राप्त्यभावात्। प्राप्तिपूर्व्वप्रतिषेधस्य  न्याय्यत्वात्। \mts यस्य सुख\-दुः\-ख\-प्राप्तिः\mte, त\vrtmm{द्विषयः सुख}{\Me(em.)}{\rdg{॰द्विषयसुख॰}{TMA}, \rdg{॰द्विषये सुख॰}{L}}दुःखा\-भाव एव योग इति सूत्रशेषप्रत्याख्यानं युज्यते॥\linelabel{sukhaduhkheend}

\llbl{p01_60}\linelabel{sasariirasya}किञ्चान्यत्---उपात्ते ऽपि सशरीर\vrtms{ग्रहणे}{\Td MA\Me}{\rdg{ग्रहणो}{\Tm}} मुक्ताद्विनिवृत्तिर्न्न \vrt{शक्यते}{\Td\LMA\Me}{\rdg{युज्यते शक्यते}{\Tm}} कर्त्तुम्, निष्प्रयोजन\-{}त्वात्। \barfu सुखदुःखाभावश्चेद्यो\mts गः विद्यमान\mte मपि शरीरं \vrtsm{कार्य्याकरणा}{\Me(em.?)}{\rdg{कार्य्या\vv{॰र्या॰}{MA}कारणा॰}{\Tm LMA}, \rdg{कार्यकारणा॰}{\Td}}दनर्त्थकं स्यात्। कार्य्यान\-{}पेक्षायां हि मुक्तामुक्तयोरविशेष एव॥\linelabel{sasariirasyaend}

\linelabel{vinigraha}%
\llbl{p01_61}किञ्चान्यत्---प्राणमनो%
\vrtmm{विनिग्रहा}{T\Me}{%
	\rdg{॰विग्रहा॰}{L\Mpc A}, 
	\rdg{॰ग्रहा॰}{\Mac}%
}%
पेक्ष इत्यपि न %
\vrt{घटते}{TMA\Me}{%
	\rdg{घ\ilg ते}{L}%
}%
, विधारकस्य %
\tstsm{प्रयत्न}{%
	See \PDhS\ \citep[70]{pdhs}: {\dn तस्य गुणाः %बुद्धि\-सुख\-दुःखे\-च्छा\-द्वेष\-%
	\ldots ॰प्रयत्न॰%
	\ldots %
	%\-धर्मा\-धर्म\-संस्कार\-संख्या\-परिमाण\-पृथक्त्व\-संयोग\-विभागाः%
	}. Cf.\ also \rVS{3.2.4}.%: {\dn %
	%प्राणा\-पान\-निमेषो\-न्मेष\-जीवन\-मनो\-गती\-न्द्रिया\-न्तरविकाराः %
%	\ldots सुख\-दुःखे इच्छा\-द्वेषौ प्रयत्न\-श्चेत्या\-त्म\-लिङ्गानि}.%
}%
स्यात्मसमवेत\-{}\varl{tvaatvi}\pause{tvaatvi}<{\dn ॰त्वात्। वि॰}>त्वात्। \tstwll{manasahkryvtvm11}{विद्यमाने ऽपि मनसः क्रियावत्त्वे}वि\donote{tvaatvi}{em.\ \rdg{॰त्वा\-द\-वि॰}{\all}}द्यमाने ऽपि म\mms न\mds सः \barfu क्रिया\vrtms{व\mte त्वे}{LM\Me}{\rdg{[\mist{1}त्वे]}{\Tm}, \rdg{\mist{1}त्व}{\Td}}\donote{manasahkryvtvm11}{Cf.\ \PDhS\ \citep[89]{pdhs}: {\dn क्रियावत्त्वान्मूर्तत्वम्}.}, % this is accidentally corect.
 \vrtsm{इन्द्रियेणा}{T\Lpc MA}{\rdg{इन्द्रि\ccs या\cce येणा}{\Lac}}संयुज्य \vrtsm{मनः}{TLM\Me}{\rdg{मतः}{A}}\vrtms{प्राणनिग्रहो}{TLMA}{\rdg{प्राण\eil मनो\eir निग्रहो}{\Me}} न कर्त्तुम्पर्य्याप्यते॥\linelabel{vinigrahaend}

\linelabel{tadanaarambha}%
\llbl{p01_62}इन्द्रिय%
\vrtms{संयोगे}{T\Lpc MA\Me}{%
	\rdg{॰संहोगे}{\Lac}%
} %
च मनसस्तदनारम्भ \mms इत्य\mme युक्तम्। %
\vrt{मनसश्च}{TL\Mpc A\Me}{%
	\rdg{मनसश्च \ccs प्रनसश्च\cce}{M}%
} %
प्राणनिग्रहे %
\vrtsm{व्यापृत}{TL\Me(em.)}{%
	\rdg{व्या\ins प्या\ine वृत॰}{M}, 
	\rdg{व्या\edl प्या\edr वृत॰}{A}, 
	\rdg{व्यापृ(वृ)त॰}{\Me}%
}% vyaap.rta is rather frequently used by Sankara
त्वादनार\-म्भा\-नु\-%
\vrtms{प\-{}पत्तिः}{TMA\Me}{%
	\rdg{॰पपतिः}{L}%
}%
। न चासतीन्द्रियवायु\-सन्निकर्षे \barfu प्रा\mms ण\mds %
\vrtms{निग्रहः}{\LMA\Me}{%
	\rdg{[\mist{3}]ा}{\Tm}, 
	\rdg{\mist{3}ा}{\Td}%
}%
\mte । न \vrtsm{च नि}{em.}{\rdg{चानि॰}{$\Sigma$}}गृहीते प्राणे योगः, आ\-त्म\-%
\vrtmm{स्थ\-विधारक}{T\Me(em.)}{%
	\rdg{॰स्वविकधारक॰}{L}, 
	\rdg{॰स्वविधारक॰}{MA}%
}%
\-प्रयत्न\-%
\vrtmm{द्वारमनो}{T}{\rdg{॰द्वारा मनो॰}{\LMA\Me}}गतक्रियासम्बन्धहेतुत्वात्प्राणनिग्रहस्य॥\linelabel{tadanaarambhaend}

\llbl{p01_63}\linelabel{vyogasamadhi}अथापि स्यात्---योगः समाधिरिति॥ 

तच्च न, \vrtsm{निष्क्रिय}{TMA\Me}{\rdg{निष्कर्य्य॰}{L}}त्वात्। आ\vrtmm{त्ममनसोर्न्नि}{em.}{\rdg{॰त्मनो नि॰}{$\Sigma$}}त्यसमाधानमेव। \tst{स्थितिश्च समाधिरित्युच्यते}{See \pageref{smdhsth},\lineref{smdhsth}.}॥% ref?
\linelabel{vyogasamadhiend}

\llbl{p01_65}तस्माद्यु\vrtmm{क्तम}{TLM\Me}{\rdg{॰क्तसम॰}{A}}भिधीयते `स च सार्व्व\vrtms{भौमश्चित्तस्य}{\Td\LMA\Me}{\rdg{॰भौमश्चित्तत्तस्य}{\Tm\ (the second {\dn त्त} beginning of a line)}} धर्म्म' इति॥


%\newpage
\llbl{p01_66}यद्येवं समाधिश्चेद्योगः, स च सार्व्वभौमश्चित्तधर्म्मः, स च सर्व्वजनानामयत्नसिद्धः, क्षिप्ताद्यनुगतत्वात्; तथा च सति यथैव श्वास\vrtmm{प्रश्वासा}{TMA\Me}{\rdg{प्रश्वसा॰}{L}}दीनाम्प्रयत्नादृते सिद्ध\vrtmm{त्वाल्लक्षणम}{TLMA}{\rdg{॰त्वात्तत्करणम\edl ल्लक्षणम\edr ॰}{\Me}}\vrtmm{नर्त्थक}{\Tmpc\Td MA\Me}{\rdg{॰नर्त्थ॰}{\Tmac L}}\-म्, एवं योगानुशासनमप्यनर्त्थकम्प्राप्तमित्यत आह---

\llbl{p01_67}\mll{atrvks1}\mula{\vrt{अत्र विक्षिप्ते}{TL}{\rdg{अनविक्षिप्ते}{MA}, \rdg{\eil तत्र\eir\  विक्षिप्ते \edl अनविक्षिप्ते\edr}{\Me}} चेतसि\labelh{VKSPTCTS1} विक्षेपोपसर्ज्जनीभूतः समाधिर्न्न योगपक्षे वर्त्तत} इति।\donote{atrvks1}{{\dn अत्र विक्षिप्ते चितसि विक्षेपोपसर्ज्जनीभूतः समाधिर्न्न योगपक्षे वर्त्तते।}} क्षिप्ताद्य\-{}नुगतस्य योगपक्षे ऽनभिप्रेतत्वात्। न हि \vrt{क्षेपाद्यनुगतः}{\Tm\Lpc MA\Me}{\rdg{विक्षेपाद्यनुगतः}{\Td}, \rdg{क्षेप\ccs त्व\cce ाद्यनुगतः}{L}} समाधिर्भूतार्त्थावद्योतनादि\vrtms{क्षमः}{TL\Mpc A\Me}{\rdg{क्षमं}{\Mac}}, क्षेपादिप्रधानत्वात्। \vrtsm{विक्षेपनि}{\LMA\Me}{\rdg{विक्षेपन्नि\vv{॰पं नि॰}{\Td}॰}{T}}\vrtms{षेधेनैव}{\Tmpc\Td\LMA\Me}{\rdg{॰षधेनैव}{\Tmac}}%\com{nsehn}{निषेधेनैव}\donote{nsehn}{\come{nsehn}}
\ हि \tstwll{pradhaanamalla}{प्रधानमल्लनिबर्हणन्यायेन}प्रधानमल्ल\vrtmm{निबर्हण}{TLM\Me}{\rdg{॰निर्बहण॰}{A}}न्याये\-न\donote{pradhaanamalla}{See \rBSBh{1.4.28}{1.4.28}; \rnBSBh{2.1.12}{2.1.12} for the use of this maxim.} क्षिप्तमूढयोरपि \barfu निषेधः \barfu कृतः। \barfu प्रधानता च विक्षेपसमाधेः, अधिकारयोग्यत्वात्। विक्षेपस्थं हि \barfu चित्तमपक्षपाता\vrtms{दिष्टं}{T\Me(em.?)}{\rdg{॰दृष्टं}{LMA}} \vrtsm{विषय}{TMA\Me}{\rdg{विषमय॰}{L}}मु\-प\-नेतुं शक्यते। न \vrt{तु}{em.}{\rdg{हि}{$\Sigma$}} \barfu विषयासञ्जनादिना \vrtsm{मूढ}{TLMA}{\rdg{क्षिप्त॰}{\Me(em.?)}}\vrtmm{मनिष्ट}{em.}{\rdg{॰मिष्ट॰}{$\Sigma$}}विषयवियोगादिना \barfu वा \vrtsm{क्षिप्त}{TLMA}{\rdg{मूढ॰}{\Me}}म\vrtmm{न्यत्रो}{TLMA}{\rdg{॰न्यथो॰}{\Me(em.?)}}पनेतुम् %\com{vryt}{वार्य्यते}%
\vrt{पार्यते}{\Me(em.)}{\rdg{[\mist{2} े]त}{\Tm}, \rdg{\mist{3}त}{\Td}, \rdg{वार्य्यते}{LMA}}%\donote{vryt}{\come{vryt}}
। `योगपक्ष' इति---समाधित्वे सत्यपि %\com{krykrn1}{कार्य्याकरणात् \ldots कार्याकरणा॰}
\vrtsm{कार्य्याकरणात्प}{\Me(em.)}{%
  \rdg{कार्य्या\vv{॰र्या॰}{M}कारणात्प॰}{\Tm LM}, %
  \rdg{कार्यकरणात्प॰}{\Td}, %
  \rdg{कार्यकारणात्प॰}{A}}क्ष इत्युच्यते। \tstwll{ngrjn}{यथा गच्छतः \ldots\ न स्थितिरित्युच्यते}यथा \vrt{गच्छतः}{TLM\Me}{\rdg{गच्छेतः}{A}} प्रतिपदं विद्यमानापि स्थितिः स्थिति%
\vrtmm{कार्य्याकरणा}{T\Me(em.)}{%
  \rdg{॰कार्य्या\vv{॰र्या॰}{MA}कारण॰}{LMA}}%\donote{krykrn1}{\come{krykrn1}}
  न्न स्थिति\-रित्युच्यते॥\donote{ngrjn}{Cf.\ \MMK\ 2.1--25.}% definitely deserves mention of Nagarjuna

%\pagebreak

\llbl{p01_68}`विक्षेपोपसर्ज्जनीभूत'त्वं विक्षेपस्थैकसमाधिव्य\ktya भिप्रायेण। `सार्व्वभौम' इति \vrt{तु}{TL\Mpc A\Me}{\rdg{{\cm om.}}{\Mac}} प्राधान्यं सामान्येन सर्व्वभूमिषु \barfu वृत्तेः॥

\llbl{p01_69}\linelabel{BhT11s}यद्युप\mts सर्ज्ज\mte नीभावः \vrtsm{योगपक्षा}{conj.}{\rdg{समाधिपक्षा॰}{$\Sigma$}}वृत्तौ कारणम्, इहाप्येकाग्र \mts एकाग्र\mte तो\vrtms{पसर्ज्जनत्वं}{\LMA\Me}{\rdg{॰पसर्ज्ज\vv{॰र्ज॰}{\Td}नं}{T}} स्यात्। ततश्च योगपक्षा\-वृत्तिरिति चेत्, अत आह---\mll{ystkgr1}\mula{यस्त्वेकाग्र} इति।
%
नैकाग्रतायाम्भूम्युपसर्ज्जनत्वम्। \barfu कुतः? 
  %\com{klskrmbl1}{\kle शकर्म्माप्रबलत्वात्। \kle शकर्म्मप्रबलत्वेन हि}%
\vrtsm{\kle शकर्म्माप्रबलत्वात्। \kle शकर्म्मप्र}{T}{
  \rdg{\ins \kle शकर्म्माप्रबलत्वात्\ine \kle शकर्म्माप्र॰}{L},
  \rdg{\kle श\-कर्मा\-प्तत्वात् \kle श\-कर्माप्र॰}{MA},
  \rdg{\kle शकर्माप्रबलत्वात्। \kle शकर्माप्र॰}{\Me}}%
बल\vrtms{त्वेन हि}{TLMA}{%
  \rdg{त्वे न हि}{\Me}}
  %\donote{klskrmbl1}{\come{klskrmbl1}} 
\barfu विक्षेपादीनामुदयः। 
\vrtsm{एका}{T\Lpc MA\Me}{%
  \rdg{एवका॰}{\Lac}%
}\-ग्रता\-भूमौ \mts स\-मा\mte धि\-र्य्यः \mula{स भूतमर्त्थं} यथात्मानम् \mula{प्रद्योतयत्य}वगमयति। 
\begin{uncertain}
अयोग्यर्त्थज्ञानमयथाभूतत्व\vrtmm{गन्धि\-त}{T\-\Lpc\-M\-A\-\Me}{%
%  \rdg{॰गन्धित॰}{T\Lpc MA\Me},
  \rdg{॰गन्थित॰}{\Lac}%
}मेवेति भूतग्रहणम्। %
\end{uncertain}\barfu
\mula{क्षिणोति} \vrt{क्षपयति}{T\Me(em.?)}{%
  \rdg{पक्षयति}{LMA}%
} \labelh{AVGRH1}\tst{पञ्चपर्व्वणो \barfu ऽविद्यादीन् \barfu \mula{\kle शान्}}{See \rYS{2.03}{2.3}: {\dn अविद्यास्मितारागद्वेषाभिनिवेशाः \barfu \kle शाः॥}.}\fubar । % ref!
\mula{कर्म्मबन्धनानि}---%\com{krmbndh1}{कर्म्माण्येव बन्धनानि}
\vrt{कर्म्माण्येव बन्धनानि}{T\Lpc \Mpc A\Me}{%
  \rdg{\om}{\Lac},
  \rdg{\ccs कर्म धर्माधर्म\cce कर्माण्येव बन्धनानि}{M}%
}%\donote{krmbndh1}{\come{krmbndh1}}
, धर्म्माधर्म्मवि%
\vrtms{मिश्राणि; यदि \barfu वा कर्म्मजानि तानि}{em.}{%
	\rdg{॰मिश्राणि यदि कर्म्मजानि तानि}{T}, 
	\rdg{॰मिश्राणि यदि कर्मजाति तानि}{\LMA}, 
	\rdg{॰विमिश्रकर्मजातानि}{\Me(em.?)}%
} %
\vrtsm{जन्म\mds ादी\mde नि ब}{\Tm\LMA}{%
	\rdg{जन्मादिनिब॰}{\Me}%
}\-न्ध\-ना\-नि---\mula{श्लथयति} शिथिलीकरोति। \mula{निरोधमामुखी\-करोत्य}भिमुखीकरोति।%
\donote{ystkgr1}{%
	{\dn यस्त्वेकाग्रे, स भूतमर्त्थम्प्रद्योतयति, क्षिणोति \kle शान्, कर्म्मबन्धनानि श्लथयति, निरोधमामुखीकरोति।}%
}%
\labelh{SSMPRJN} %
\mll{SSMPRJ1}\mula{स सम्प्रज्ञात \barfu इत्या}%
\vrtms{\mula{ख्यायत}}{T}{%
	\rdg{॰ख्यायते}{\LMA\Me}%
}%
\donote{SSMPRJ1}{%
	{\dn  स सम्प्रज्ञात इत्याख्यायते।}%
} \barfu आचार्य्यैः॥\linelabel{BhT11e} % This aacaarya must be the author of the YBh. -> YVi-kaara uses plural to refer to a single author.

\llbl{p01_70}\mll{vtrknugt1}\mula{स च वितर्क्कानु\vrtmm{गतो विचारा}{\Td \Me(em.)}{\rdg{॰गतो पि\vv{॰गतो ऽपि}{MA} विचारा॰}{TLMA}}नुगत} इत्यादि भाष्यम्---सम्प्र\vrtmm{ज्ञातस्य प्रा}{TL\Mpc A\Me}{\rdg{॰ज्ञातस्य \ccs सा\cce प्रा॰}{M}}थम्यात्। लक्ष्णा\-{}भिधाने \vrt{तु}{conj.}{\rdg{\om}{$\Sigma$}} प्राप्ते, \barfu अभ्यर्हितत्वादसम्प्रज्ञा\vrtms{त\-\mds स्यै\mde\mms व लक्ष\mme णं}{\LMA\Me}{\rdg{॰तस्ये[\mist{3}]णं}{\Tm}, \rdg{॰त\mist{1}पलक्षणं}{\Td}} युक्तमिह वक्तुम्---\mula{इति} \tstwll{refys117}{उपरिष्टात्प्रवेदयिष्यामः}सम्प्रज्ञातमु\-\mula{प\-रिष्टा\-त्प्रवेदयिष्याम}\donote{refys117}{See \rYS{1.17}{1.17}: {\dn वितर्कविचारानन्दास्मितारूपानुगमात् \barfu संप्रज्ञातः॥}.} इत्युच्यते॥\donote{vtrknugt1}{{\dn स च \barfu वितर्क्कानुगतो विचारानुगत \barfu इत्युपरिष्टात्प्रवेदयिष्यामः।}}%
\fakepar
\llbl{p01_71}इतश्चासम्प्रज्ञातस्य लक्षणमिहैव व\mts क्तव्यम््---सम्प्र\mte ज्ञातनिरपेक्षो \barfu ऽप्यसम्प्रज्ञातः \barfu  \vrtsm{प्रकृष्ट}{TLMA}{\rdg{पर\edl प्रकृष्ट\edr ॰}{\Me}}\-वै\-रा\-ग्य\-विरामप्रत्यया\vrtms{भ्यासाभ्यां}{\Td\LMA\Me}{\rdg{॰भ्यासाभ्यासाभ्यां}{\Tm}} सिध्यतीत्येतत्प्रदर्श\mms ना\mme र्त्थम्। इह \vrt{तु}{T}{\rdg{त}{L}, \rdg{\om}{MA\Me}} सम्प्रज्ञातलक्षणाभिधाने तदुत्तरकाले चासम्प्रज्ञातलक्ष\vrtmm{णाभिधाने सम्प्र}{T\-\Me(em.)}{\rdg{॰णमभिधानम्प्र\vv{॰नं प्र॰}{MA}॰}{LMA}}\labelh{TEcorr2}\mts ज्ञातापेक्षयै\vrtmm{वासम्प्र}{T}{\rdg{॰व असम्प्र\vv{॰संप्र॰}{MA}॰}{\LMA\Me}}\-ज्ञात\vrtmm{स\-मा\-धा\mde व\-धि}{LM\Me}{\rdg{॰समानाधि॰}{A}}\mme कार \barfu इत्याशंका स्या\-त्। तस्मादु`परिष्टात्प्रवेदयिष्याम' इत्याह॥१॥
